\topic{进阶综合题卡}

\noindent\textbf{LeetCode 4 寻找两个正序数组的中位数。}

\textbf{题目说明。} LeetCode 4 寻找两个正序数组的中位数 的题面入口要先还原成具体任务：两个有序数组的中位数可以看成左右两半切分后的边界值。

\textbf{样例入口。} 拿 [1, 3] 和 [2] 手算，合并后中位数是 2；但不用真合并，只要把总长度切成左半两个数、右半一个数。再拿 [1, 2] 和 [3, 4]，如果短数组切 1 个、长数组切 1 个，左半最大 3 大于右半最小 2，切分错了，应把短数组切分往右调。

\textbf{暴力基线。} 慢版本按顺序扫描下标、逐个试答案值，或直接检查候选直到遇到目标边界。它先保证候选空间覆盖完整，但每次只排除一个候选，浪费了题目给出的顺序、比较或可行性边界。放到本题就是：先按“两个有序数组的中位数可以看成左右两半切分后的边界值”完整试慢版本；样例里已经能看到“规律是二分短数组切分点，使左半元素个数正确且左右边界有序”，所以优化前必须先能说清慢版本枚举了哪些候选。

\textbf{暴力过程。} 拿 [1, 3] 和 [2] 手算，合并后中位数是 2；但不用真合并，只要把总长度切成左半两个数、右半一个数。再拿 [1, 2] 和 [3, 4]，如果短数组切 1 个、长数组切 1 个，左半最大 3 大于右半最小 2，切分错了，应把短数组切分往右调。

\textbf{重复工作。} 题面模型：两个有序数组的中位数可以看成左右两半切分后的边界值。暴力版的慢点要落到本题：慢版本会按顺序试“两个有序数组的中位数可以看成左右两半切分后的边界值”里的候选；而样例规律已经指出“规律是二分短数组切分点，使左半元素个数正确且左右边界有序”。一次比较或一次可行性检查能丢掉一整段候选，继续线性试就是浪费。后面的优化只消掉这类重复工作，不能改变答案集合。

\textbf{优化条件。} 本题的关键观察是：二分较短数组的切分位置，让左半元素数量正确且左半最大值不大于右半最小值。这句话必须能翻译成可维护状态；如果题目条件改变后观察不再成立，就不能继续套原来的写法。

\textbf{相邻状态。} 规律是二分短数组切分点，使左半元素个数正确且左右边界有序。把这个特例继续拆开看：每次比较 mid 之后，必须能指出哪一侧整体不可能包含目标，或哪一侧整体已经满足可行性。二分的核心不是 mid 公式，而是被丢弃区间确实不含答案。

\textbf{状态复用。} 把逐个试候选改成维护答案所在区间；边界谓词、可行性检查或局部比较给出分支方向，每个分支都必须能排除一侧候选。本题落点是：规律是二分短数组切分点，使左半元素个数正确且左右边界有序。手算时只改被新输入影响的那一小部分，而不是回头重算完整候选。

\textbf{指针和字段。} 状态字段要能说清：left、right、mid、mid 处比较值或检查返回值、答案区间含义、返回值语义；每一轮更新都要保留目标位置或第一个可行值。手算时按这个协议跟踪：拿 [1, 3] 和 [2] 手算，合并后中位数是 2；但不用真合并，只要把总长度切成左半两个数、右半一个数。再拿 [1, 2] 和 [3, 4]，如果短数组切 1 个、长数组切 1 个，左半最大 3 大于右半最小 2，切分错了，应把短数组切分往右调。然后按协议复查：手算时列出 left、mid、right，以及 mid 处比较结果、可行性检查结果或局部趋势。每一步都要写出被丢弃的一侧为什么不含答案。

\textbf{代码落点。} 检查函数或比较分支单独写清，mid 不溢出，循环退出条件和返回值配套；每个分支都要解释丢弃区间。关键代码行必须能逐句翻译成“旧状态怎样变成新状态”，否则就只是模板。

\textbf{正确性。} 证明从排除一侧开始：答案空间题证明谓词单调，局部比较题证明保留下来的区间仍含目标；不能只靠经验猜测左右边界。

\textbf{边界实验。} 先用某个数组为空、总长度奇偶、切分在边界、哨兵值不能参与真实比较做反例检查。实验时覆盖某个数组为空、总长度奇偶、切分在边界、哨兵值不能参与真实比较，逐轮打印 left、mid、right、比较值或检查结果。只要有一次被排除区间仍含答案，二分不变量就错了。

\textbf{复杂度变化。} 暴力逐个试候选是线性或和值域成正比；二分每次排除一半候选，复杂度变成对数级，答案二分还要乘上单次检查成本。

\textbf{迁移追问。} 如果题目改成“若只要求第 k 小元素，可以用递归丢弃较小前缀或二分切分的同类思想”，先判断原状态是否仍然覆盖新答案集合，再决定是微调边界、改 value 语义，还是换算法。

\noindent\textbf{LeetCode 10 正则表达式匹配。}

\textbf{题目说明。} LeetCode 10 正则表达式匹配 的题面入口要先还原成具体任务：模式中的点号匹配单字符，星号修饰前一个字符并允许出现零次或多次。

\textbf{样例入口。} 拿 s=a、p=a* 手算，星号可以让 a 出现一次；拿 s=空、p=a*，星号又可以让 a 出现零次。再拿 s=aa、p=a*，消耗一个 a 后模式仍停在 a*，还能继续匹配。

\textbf{暴力基线。} 慢版本先写递归搜索树：节点表示处理位置、剩余资源和当前选择。只要不同路径反复到达同一个节点，就说明需要把这个节点命名成 DP 状态。放到本题就是：先按“模式中的点号匹配单字符，星号修饰前一个字符并允许出现零次或多次”完整试慢版本；样例里已经能看到“规律是星号有两条分叉：当零次用时看 dp[i][j-2]，当消耗当前字符时看 dp[i-1][j]”，所以优化前必须先能说清慢版本枚举了哪些候选。

\textbf{暴力过程。} 拿 s=a、p=a* 手算，星号可以让 a 出现一次；拿 s=空、p=a*，星号又可以让 a 出现零次。再拿 s=aa、p=a*，消耗一个 a 后模式仍停在 a*，还能继续匹配。

\textbf{重复工作。} 题面模型：模式中的点号匹配单字符，星号修饰前一个字符并允许出现零次或多次。暴力版的慢点要落到本题：慢版本会在“模式中的点号匹配单字符，星号修饰前一个字符并允许出现零次或多次”的选择树里反复到达相同参数。样例规律是“规律是星号有两条分叉：当零次用时看 dp[i][j-2]，当消耗当前字符时看 dp[i-1][j]”，说明这个参数组合可以命名成状态，算过之后后面直接复用。后面的优化只消掉这类重复工作，不能改变答案集合。

\textbf{优化条件。} 本题的关键观察是：二维 DP 表示两个前缀是否匹配，星号转移同时覆盖零次和消耗一个字符后继续匹配。这句话必须能翻译成可维护状态；如果题目条件改变后观察不再成立，就不能继续套原来的写法。

\textbf{相邻状态。} 规律是星号有两条分叉：当零次用时看 dp[i][j-2]，当消耗当前字符时看 dp[i-1][j]。把这个特例继续拆开看：先问暴力搜索里哪些不同路径到达同一个后续问题，再给这个后续问题命名为状态。状态必须包含未来决策需要的全部信息，转移只从更小且已经算清的状态来。

\textbf{状态复用。} 把重复递归节点命名为状态；遍历顺序只是在保证依赖状态已经算完。本题落点是：规律是星号有两条分叉：当零次用时看 dp[i][j-2]，当消耗当前字符时看 dp[i-1][j]。手算时只改被新输入影响的那一小部分，而不是回头重算完整候选。

\textbf{指针和字段。} 状态字段要能说清：状态下标、状态值语义、初始化、转移来源、遍历顺序；空间压缩时还要指出读到的是旧值还是新值。手算时按这个协议跟踪：拿 s=a、p=a* 手算，星号可以让 a 出现一次；拿 s=空、p=a*，星号又可以让 a 出现零次。再拿 s=aa、p=a*，消耗一个 a 后模式仍停在 a*，还能继续匹配。然后按协议复查：手算时先写一个很小的 DP 表或记忆化搜索记录。每个格子旁边写它来自哪些更小状态。

\textbf{代码落点。} 先写未压缩版本校验转移；压缩前后用小表对拍；不可达哨兵参与加法前要检查溢出。关键代码行必须能逐句翻译成“旧状态怎样变成新状态”，否则就只是模板。

\textbf{正确性。} 证明从最后一步开始：任意最优解的最后一步都在转移枚举中，转移来源已经由更小状态正确计算。

\textbf{边界实验。} 先用空串、能匹配空串的星号链、星号前必须有字符、点号和普通字符的匹配条件做反例检查。实验时覆盖空串、能匹配空串的星号链、星号前必须有字符、点号和普通字符的匹配条件，先打印完整 DP 表，再比较空间压缩版本。若压缩版不同，优先检查遍历顺序和旧值保存。

\textbf{复杂度变化。} 暴力递归会重复计算相同子问题，常见指数级；DP 把每个状态算一次，时间是状态数乘单个转移成本，空间是状态表大小或压缩后的大小。

\textbf{迁移追问。} 如果题目改成“通配符匹配中的星号语义是任意长度字符串，不能直接搬用本题转移”，先判断原状态是否仍然覆盖新答案集合，再决定是微调边界、改 value 语义，还是换算法。

\noindent\textbf{LeetCode 23 合并 K 个升序链表。}

\textbf{题目说明。} LeetCode 23 合并 K 个升序链表 的题面入口要先还原成具体任务：每条链表当前头节点都是可能的下一个最小元素。

\textbf{样例入口。} 拿三条链表头 1、1、2 手算，每次答案只需要当前最小头节点；弹出某条链的头后，只有它的后继会成为新候选。暴力每轮扫 k 个头，堆把“找当前最小头”变成对数操作。

\textbf{暴力基线。} 暴力每轮扫描 K 个链表头，找到当前最小头节点接到结果后推进该链。

\textbf{暴力过程。} 若三条链当前头是 1、1、2，第一步只需从这三个头里选最小；弹出某条链的 1 后，只有它的后继会成为新候选。

\textbf{重复工作。} 题面模型：每条链表当前头节点都是可能的下一个最小元素。暴力版的慢点要落到本题：浪费在每输出一个节点都线性扫 K 个头。候选集合动态变化，但每次只需要最小头节点，适合小顶堆。后面的优化只消掉这类重复工作，不能改变答案集合。

\textbf{优化条件。} 本题的关键观察是：小顶堆保存每条非空链的当前头，弹出最小节点后把它的后继入堆。这句话必须能翻译成可维护状态；如果题目条件改变后观察不再成立，就不能继续套原来的写法。

\textbf{相邻状态。} 规律是堆里始终最多保存每条非空链的一个当前头。把这个特例继续拆开看：堆顶必须正好代表下一步最需要处理的候选；若堆里可能有过期对象，读取堆顶前要先清理。堆只解决动态极值，不自动证明被弹出或被丢弃的元素无用。

\textbf{状态复用。} 小顶堆保存每条非空链的当前头节点。弹出最小节点接入结果；若该节点有 next，把 next 入堆。手算时只改被新输入影响的那一小部分，而不是回头重算完整候选。

\textbf{指针和字段。} 状态字段要能说清：heap 中元素是 ListNode*，比较依据是 node->val；tail 是结果尾；node->next 是该链的新候选。手算时按这个协议跟踪：头节点 1、1、2 入堆；弹出第一个 1 后把它的后继入堆，堆仍保存每条链当前最小暴露节点。

\textbf{代码落点。} 关键是堆里保存节点指针而不是复制值，这样可以复用原节点。比较器要写成小顶堆语义，C++ priority\_queue 默认是大顶堆。关键代码行必须能逐句翻译成“旧状态怎样变成新状态”，否则就只是模板。

\textbf{正确性。} 每条有序链未输出部分的最小值就是当前头；所有链当前头中的最小者就是全局下一个输出节点。

\textbf{边界实验。} 先用空链表数组、某条链为空、重复值、堆里保存节点指针还是值做反例检查。实验时覆盖空链表数组、某条链为空、重复值、堆里保存节点指针还是值，记录堆顶是否过期、比较器是否让正确候选在堆顶、K 或窗口大小为边界值时是否稳定。

\textbf{复杂度变化。} 每轮扫 K 个头是 O(NK)。堆中最多 K 个节点，每个节点入堆出堆一次，总时间 O(N log K)，空间 O(K)。

\textbf{迁移追问。} 如果题目改成“也可以分治两两合并，堆更适合 K 较大且链长不均的场景”，先判断原状态是否仍然覆盖新答案集合，再决定是微调边界、改 value 语义，还是换算法。

\noindent\textbf{LeetCode 25 K 个一组翻转链表。}

\textbf{题目说明。} LeetCode 25 K 个一组翻转链表 的题面入口要先还原成具体任务：链表被切成长度为 K 的连续组，只有完整组才反转。

\textbf{样例入口。} 拿 1->2->3->4、k=2 画图。第一组边界是 1 和 2，下一组入口是 3；反转前如果不保存 3，链会断。反转后上一组尾接 2，旧头 1 接 3。

\textbf{暴力基线。} 暴力可以把每 K 个节点的值放进数组反转，再写回链表；这避开了指针重连，但没有真正按节点反转。

\textbf{暴力过程。} 以 1->2->3->4->5、k=2 为例，第一组 1、2 反转成 2->1，第二组 3、4 反转成 4->3，剩余 5 不动。

\textbf{重复工作。} 题面模型：链表被切成长度为 K 的连续组，只有完整组才反转。暴力版的慢点要落到本题：浪费在数组辅助和节点值交换。题目要求链表节点重排时，应直接反转每一组的 next 指针。后面的优化只消掉这类重复工作，不能改变答案集合。

\textbf{优化条件。} 本题的关键观察是：每轮先探测 K 个节点确认完整，再保存下一组头并做局部反转接回。这句话必须能翻译成可维护状态；如果题目条件改变后观察不再成立，就不能继续套原来的写法。

\textbf{相邻状态。} 规律是每组先找 kth 和 group\_next，再局部反转并接回前后两段。把这个特例继续拆开看：每个指针都要有角色，上一段尾、当前段头、当前节点和后继入口不能混。只要一次改 next 之前没有保存后继，就可能丢掉整段未处理链。

\textbf{状态复用。} dummy 指向头。group\_prev 是上一组尾，先从 group\_prev 出发找第 k 个节点 group\_end；不足 k 个就停止。保存 next\_group 后反转 [group\_start, group\_end]，再接回。手算时只改被新输入影响的那一小部分，而不是回头重算完整候选。

\textbf{指针和字段。} 状态字段要能说清：group\_prev 是已处理部分尾，group\_start 是本组头，group\_end 是本组尾，next\_group 是下一组入口。手算时按这个协议跟踪：第一组 start=1、end=2、next\_group=3。反转后得到 2->1，再让 group\_prev->next=2，1->next=3。

\textbf{代码落点。} 关键顺序是先保存 next\_group，再局部反转；反转结束后原 group\_start 变成本组尾，下一轮 group\_prev=group\_start。关键代码行必须能逐句翻译成“旧状态怎样变成新状态”，否则就只是模板。

\textbf{正确性。} 每次只改变完整 K 组内部指针，并把本组头尾接回已处理链和未处理链；不足 K 个节点不进入反转，符合题意。

\textbf{边界实验。} 先用不足 K 个不反转、K 为一、刚好整除、反转后头尾接错做反例检查。实验时覆盖不足 K 个不反转、K 为一、刚好整除、反转后头尾接错，每步画出前驱、当前、后继和下一段头。链表题不要只看输出值，还要检查节点是否丢失或成环。

\textbf{复杂度变化。} 数组辅助 O(n) 空间；原地分组反转每节点访问常数次，O(n) 时间、O(1) 空间。

\textbf{迁移追问。} 如果题目改成“递归写法短但可能有栈风险，迭代写法更接近工程实现”，先判断原状态是否仍然覆盖新答案集合，再决定是微调边界、改 value 语义，还是换算法。

\noindent\textbf{LeetCode 32 最长有效括号。}

\textbf{题目说明。} LeetCode 32 最长有效括号 的题面入口要先还原成具体任务：有效括号子串需要知道每段非法边界之后的最早可连接位置。

\textbf{样例入口。} 拿字符串 )()()) 手算。第一个 ) 不是可匹配对象，而是非法边界；后面每次匹配成功，当前有效长度由当前位置减去栈顶边界得出。若栈空，要把当前右括号作为新非法边界压入。

\textbf{暴力基线。} 慢版本让每个位置向左或向右线性寻找边界，或者让每个结构反复等待匹配。它能算出答案，但很多尚未完成的候选会被未来同一个事件一起解决。放到本题就是：先按“有效括号子串需要知道每段非法边界之后的最早可连接位置”完整试慢版本；样例里已经能看到“规律是栈里既保存未匹配左括号下标，也保存最近非法边界”，所以优化前必须先能说清慢版本枚举了哪些候选。

\textbf{暴力过程。} 拿字符串 )()()) 手算。第一个 ) 不是可匹配对象，而是非法边界；后面每次匹配成功，当前有效长度由当前位置减去栈顶边界得出。若栈空，要把当前右括号作为新非法边界压入。

\textbf{重复工作。} 题面模型：有效括号子串需要知道每段非法边界之后的最早可连接位置。暴力版的慢点要落到本题：慢版本会围绕“有效括号子串需要知道每段非法边界之后的最早可连接位置”反复寻找最近边界或最近未完成对象。样例规律是“规律是栈里既保存未匹配左括号下标，也保存最近非法边界”，说明旧候选一旦遇到当前输入就能被批量结算；继续为每个位置单独向左或向右扫描就是浪费。后面的优化只消掉这类重复工作，不能改变答案集合。

\textbf{优化条件。} 本题的关键观察是：栈保存下标，先放入哨兵负一；匹配后用当前下标减栈顶得到长度。这句话必须能翻译成可维护状态；如果题目条件改变后观察不再成立，就不能继续套原来的写法。

\textbf{相邻状态。} 规律是栈里既保存未匹配左括号下标，也保存最近非法边界。把这个特例继续拆开看：栈里的对象都是答案尚未确定的候选；一旦当前元素触发弹栈，被弹出的候选必须已经同时拿到了左边界和右边界，未来元素不再有资格改变它的答案。

\textbf{状态复用。} 把反复寻找最近边界改成维护未完成候选；新元素只和栈顶比较，连续弹出一批已经被当前元素解决的对象。本题落点是：规律是栈里既保存未匹配左括号下标，也保存最近非法边界。手算时只改被新输入影响的那一小部分，而不是回头重算完整候选。

\textbf{指针和字段。} 状态字段要能说清：栈内对象的业务含义、当前输入、弹出时结算的对象、左边界和右边界；栈中存下标还是值要明确。手算时按这个协议跟踪：拿字符串 )()()) 手算。第一个 ) 不是可匹配对象，而是非法边界；后面每次匹配成功，当前有效长度由当前位置减去栈顶边界得出。若栈空，要把当前右括号作为新非法边界压入。然后按协议复查：手算时记录栈内元素的业务含义，不只记录数值。入栈表示答案未定，弹栈表示当前事件已经让它的答案定下来。

\textbf{代码落点。} 弹栈前确认栈非空，弹出后再读取新的左边界；相等元素和哨兵高度要有统一策略。关键代码行必须能逐句翻译成“旧状态怎样变成新状态”，否则就只是模板。

\textbf{正确性。} 证明从弹出时机开始：被弹出的对象在当前时刻答案已经确定，未来元素无法再改变它。

\textbf{边界实验。} 先用空串、全左括号、全右括号、多个断开的合法段做反例检查。实验时准备空串、全左括号、全右括号、多个断开的合法段，打印每次入栈、弹栈和结算对象。重点观察相等元素、空栈、哨兵和最后一次结算。

\textbf{复杂度变化。} 暴力为每个位置寻找边界常见为平方级；单调栈让每个元素最多入栈一次、出栈一次，因此主体扫描为线性时间。

\textbf{迁移追问。} 如果题目改成“DP 写法也可做，但栈写法更直接表达最近非法边界”，先判断原状态是否仍然覆盖新答案集合，再决定是微调边界、改 value 语义，还是换算法。

\noindent\textbf{LeetCode 41 缺失的第一个正数。}

\textbf{题目说明。} LeetCode 41 缺失的第一个正数 的题面入口要先还原成具体任务：答案只可能落在一到 n 加一之间，数组本身可以充当值到位置的映射。

\textbf{样例入口。} 拿 [3, 4, -1, 1] 手算。答案只可能在 1 到 n+1；值 3 应放到下标 2，值 4 放到下标 3，-1 不管，1 放到下标 0。放完后第一个位置不等于 i+1 的地方就是缺失值。再试 [1, 1]，若不判断目标位置已有相同值会无限交换。

\textbf{暴力基线。} 暴力可以从 1 开始逐个查找是否出现在数组中，或者用哈希集合记录所有正数后再查。

\textbf{暴力过程。} 以 [3, 4, -1, 1] 为例，答案只可能在 1..5。把 3 放到下标 2，把 4 放到下标 3，把 1 放到下标 0，最后下标 1 缺 2。

\textbf{重复工作。} 题面模型：答案只可能落在一到 n 加一之间，数组本身可以充当值到位置的映射。暴力版的慢点要落到本题：浪费在额外集合或重复查找。长度为 n 的数组只关心 1..n 这些值，数组位置本身可以当哈希桶。后面的优化只消掉这类重复工作，不能改变答案集合。

\textbf{优化条件。} 本题的关键观察是：把合法正数 x 尽量交换到下标 x 减一，最后第一个位置不匹配就是答案。这句话必须能翻译成可维护状态；如果题目条件改变后观察不再成立，就不能继续套原来的写法。

\textbf{相邻状态。} 规律是数组本身当哈希表，只处理合法正数且避免重复交换。把这个特例继续拆开看：先标出已经确认的前缀或后缀，再标出未知区；能被丢掉的候选，必须是以后再也不会改变已确认区域语义的候选。这样才算从例子推出数组不变量，而不是只记一次操作。

\textbf{状态复用。} 扫描每个位置，把合法值 x 交换到 nums[x-1]。若目标位置已经是 x，停止交换，避免重复值死循环。手算时只改被新输入影响的那一小部分，而不是回头重算完整候选。

\textbf{指针和字段。} 状态字段要能说清：i 是当前整理位置，x=nums[i] 是试图归位的值，x-1 是它的目标下标。手算时按这个协议跟踪：[3, 4, -1, 1] 中 i=0，3 应去下标 2，换来 -1；i=1，4 应去下标 3，换来 1；1 再换到下标 0。

\textbf{代码落点。} 关键 while 条件是 1<=nums[i]<=n \&\& nums[nums[i]-1]!=nums[i]。第二遍找第一个 nums[i]!=i+1。关键代码行必须能逐句翻译成“旧状态怎样变成新状态”，否则就只是模板。

\textbf{正确性。} 每次有效交换都会至少让一个合法正数回到正确位置；最终若位置 i 不等于 i+1，则 i+1 没有被任何位置放回来，就是最小缺失正数。

\textbf{边界实验。} 先用非正数、大于 n 的数、重复值导致无限交换、数组已经完整包含一到 n做反例检查。实验时把非正数、大于 n 的数、重复值导致无限交换、数组已经完整包含一到 n列成表，再打印关键下标和有效区间。若某个元素被写入后又被未来步骤破坏，说明区间不变量没有成立。

\textbf{复杂度变化。} 哈希法 O(n) 空间；原地归位每个数交换有限次，总时间 O(n)，额外空间 O(1)。

\textbf{迁移追问。} 如果题目改成“若不能修改输入，只能使用哈希集合或额外布尔数组换取更简单边界”，先判断原状态是否仍然覆盖新答案集合，再决定是微调边界、改 value 语义，还是换算法。

\noindent\textbf{LeetCode 44 通配符匹配。}

\textbf{题目说明。} LeetCode 44 通配符匹配 的题面入口要先还原成具体任务：问号匹配一个字符，星号匹配任意长度字符串，状态是两个前缀能否匹配。

\textbf{样例入口。} 拿 s=ab、p=a* 手算，星号可以吞掉 b；拿 s=空、p=*，星号也可以代表空。DP 中星号的两条路是匹配空串或继续吞当前字符。贪心中失配时回到最近星号，让它多吞一个字符。

\textbf{暴力基线。} 慢版本先写递归搜索树：节点表示处理位置、剩余资源和当前选择。只要不同路径反复到达同一个节点，就说明需要把这个节点命名成 DP 状态。放到本题就是：先按“问号匹配一个字符，星号匹配任意长度字符串，状态是两个前缀能否匹配”完整试慢版本；样例里已经能看到“规律是本题星号独立代表任意串，和正则里修饰前一字符的星号不同”，所以优化前必须先能说清慢版本枚举了哪些候选。

\textbf{暴力过程。} 拿 s=ab、p=a* 手算，星号可以吞掉 b；拿 s=空、p=*，星号也可以代表空。DP 中星号的两条路是匹配空串或继续吞当前字符。贪心中失配时回到最近星号，让它多吞一个字符。

\textbf{重复工作。} 题面模型：问号匹配一个字符，星号匹配任意长度字符串，状态是两个前缀能否匹配。暴力版的慢点要落到本题：慢版本会在“问号匹配一个字符，星号匹配任意长度字符串，状态是两个前缀能否匹配”的选择树里反复到达相同参数。样例规律是“规律是本题星号独立代表任意串，和正则里修饰前一字符的星号不同”，说明这个参数组合可以命名成状态，算过之后后面直接复用。后面的优化只消掉这类重复工作，不能改变答案集合。

\textbf{优化条件。} 本题的关键观察是：DP 中星号可匹配空串或继续吞掉当前字符；贪心写法记录最近星号并在失配时回退扩展星号覆盖范围。这句话必须能翻译成可维护状态；如果题目条件改变后观察不再成立，就不能继续套原来的写法。

\textbf{相邻状态。} 规律是本题星号独立代表任意串，和正则里修饰前一字符的星号不同。把这个特例继续拆开看：先问暴力搜索里哪些不同路径到达同一个后续问题，再给这个后续问题命名为状态。状态必须包含未来决策需要的全部信息，转移只从更小且已经算清的状态来。

\textbf{状态复用。} 把重复递归节点命名为状态；遍历顺序只是在保证依赖状态已经算完。本题落点是：规律是本题星号独立代表任意串，和正则里修饰前一字符的星号不同。手算时只改被新输入影响的那一小部分，而不是回头重算完整候选。

\textbf{指针和字段。} 状态字段要能说清：状态下标、状态值语义、初始化、转移来源、遍历顺序；空间压缩时还要指出读到的是旧值还是新值。手算时按这个协议跟踪：拿 s=ab、p=a* 手算，星号可以吞掉 b；拿 s=空、p=*，星号也可以代表空。DP 中星号的两条路是匹配空串或继续吞当前字符。贪心中失配时回到最近星号，让它多吞一个字符。然后按协议复查：手算时先写一个很小的 DP 表或记忆化搜索记录。每个格子旁边写它来自哪些更小状态。

\textbf{代码落点。} 先写未压缩版本校验转移；压缩前后用小表对拍；不可达哨兵参与加法前要检查溢出。关键代码行必须能逐句翻译成“旧状态怎样变成新状态”，否则就只是模板。

\textbf{正确性。} 证明从最后一步开始：任意最优解的最后一步都在转移枚举中，转移来源已经由更小状态正确计算。

\textbf{边界实验。} 先用连续星号、空串、模式只含星号、贪心回退时字符串位置不能倒退错做反例检查。实验时覆盖连续星号、空串、模式只含星号、贪心回退时字符串位置不能倒退错，先打印完整 DP 表，再比较空间压缩版本。若压缩版不同，优先检查遍历顺序和旧值保存。

\textbf{复杂度变化。} 暴力递归会重复计算相同子问题，常见指数级；DP 把每个状态算一次，时间是状态数乘单个转移成本，空间是状态表大小或压缩后的大小。

\textbf{迁移追问。} 如果题目改成“正则表达式匹配的星号修饰前一个字符，和本题星号独立匹配任意串不同”，先判断原状态是否仍然覆盖新答案集合，再决定是微调边界、改 value 语义，还是换算法。

\noindent\textbf{LeetCode 51 N 皇后。}

\textbf{题目说明。} LeetCode 51 N 皇后 的题面入口要先还原成具体任务：每一行放一个皇后，列和两条对角线不能冲突。

\textbf{样例入口。} 拿 n=4 手算，第一行放某列后，下一行不能放同列、主对角线和副对角线冲突的位置。

\textbf{暴力基线。} 慢版本就是完整搜索树：每层列出选择列表，路径到达终止条件时检查答案。它先保证覆盖所有可能，再把明显无效或重复的分支剪掉。放到本题就是：先按“每一行放一个皇后，列和两条对角线不能冲突”完整试慢版本；样例里已经能看到“规律是按行递归，每行放一个皇后，三个集合分别记录列、row-col、row+col 是否被占用”，所以优化前必须先能说清慢版本枚举了哪些候选。

\textbf{暴力过程。} 拿 n=4 手算，第一行放某列后，下一行不能放同列、主对角线和副对角线冲突的位置。

\textbf{重复工作。} 题面模型：每一行放一个皇后，列和两条对角线不能冲突。暴力版的慢点要落到本题：慢版本会围绕“每一行放一个皇后，列和两条对角线不能冲突”展开完整搜索树。样例规律是“规律是按行递归，每行放一个皇后，三个集合分别记录列、row-col、row+col 是否被占用”，说明一层递归必须对应一个明确决策，剪枝只能删除整棵不可能成功或重复的子树。后面的优化只消掉这类重复工作，不能改变答案集合。

\textbf{优化条件。} 本题的关键观察是：逐行搜索，用列集合、主对角线集合和副对角线集合快速判断合法性。这句话必须能翻译成可维护状态；如果题目条件改变后观察不再成立，就不能继续套原来的写法。

\textbf{相邻状态。} 规律是按行递归，每行放一个皇后，三个集合分别记录列、row-col、row+col 是否被占用。把这个特例继续拆开看：一层递归对应一个明确决策，路径保存已经做出的选择，选择列表保存仍可能扩展的分支。剪枝前必须能说明被剪掉的整棵子树没有合法答案，而不是只因为它看起来不优。

\textbf{状态复用。} 把完整搜索树加上合法性检查、剩余资源剪枝和同层去重；路径状态进入和退出要成对恢复。本题落点是：规律是按行递归，每行放一个皇后，三个集合分别记录列、row-col、row+col 是否被占用。手算时只改被新输入影响的那一小部分，而不是回头重算完整候选。

\textbf{指针和字段。} 状态字段要能说清：路径、起点或使用标记、剩余目标、答案集合、剪枝条件；进入递归和退出递归必须成对恢复。手算时按这个协议跟踪：拿 n=4 手算，第一行放某列后，下一行不能放同列、主对角线和副对角线冲突的位置。然后按协议复查：手算时给每层标出选择列表和路径。剪枝发生前要指出被剪分支为什么已经不可能得到合法答案。

\textbf{代码落点。} 剪枝条件写在扩展前；同层去重不要误写成全局去重；保存答案时复制当前路径。关键代码行必须能逐句翻译成“旧状态怎样变成新状态”，否则就只是模板。

\textbf{正确性。} 证明从搜索树覆盖开始：每个合法答案有且只有一条路径，剪枝只删掉不可能成功的分支。

\textbf{边界实验。} 先用n 为一、n 为二或三无解、对角线编号、回退时状态恢复做反例检查。实验时覆盖n 为一、n 为二或三无解、对角线编号、回退时状态恢复，统计搜索节点数量和剪枝次数。重复元素题要打印同一层跳过哪些值，确认没有误删合法路径。

\textbf{复杂度变化。} 暴力搜索树的规模由输出或选择空间决定；剪枝不改变最坏指数级本质，但能删除不可能成功或重复的分支，复杂度要说明输出规模。

\textbf{迁移追问。} 如果题目改成“可用位掩码压缩三个约束集合，提高常数性能”，先判断原状态是否仍然覆盖新答案集合，再决定是微调边界、改 value 语义，还是换算法。

\noindent\textbf{LeetCode 72 编辑距离。}

\textbf{题目说明。} LeetCode 72 编辑距离 的题面入口要先还原成具体任务：二维状态表示两个前缀之间的最少编辑操作。

\textbf{样例入口。} 拿 word1=ab、word2=ac 手算最后一个字符。若 b 改成 c，是 dp[1][1]+1；若删掉 b，是 dp[1][2]+1；若插入 c，是 dp[2][1]+1。字符相等时才直接继承左上。

\textbf{暴力基线。} 暴力递归比较两个字符串后缀：相等就同时跳过，不等就尝试插入、删除、替换三种操作。

\textbf{暴力过程。} 以 word1=ab、word2=ac 为例，末尾 b 和 c 不同。替换 b 成 c 来自 dp[1][1]+1；删除 b 来自 dp[1][2]+1；插入 c 来自 dp[2][1]+1。

\textbf{重复工作。} 题面模型：二维状态表示两个前缀之间的最少编辑操作。暴力版的慢点要落到本题：浪费在不同编辑序列会反复到达同一对前缀长度。只要 i 和 j 相同，后续最少编辑次数就相同。后面的优化只消掉这类重复工作，不能改变答案集合。

\textbf{优化条件。} 本题的关键观察是：末尾字符相同继承左上；不同则枚举插入、删除、替换。这句话必须能翻译成可维护状态；如果题目条件改变后观察不再成立，就不能继续套原来的写法。

\textbf{相邻状态。} 规律是每个格子枚举最后一次编辑操作，而不是凭感觉填三邻居最小。把这个特例继续拆开看：先问暴力搜索里哪些不同路径到达同一个后续问题，再给这个后续问题命名为状态。状态必须包含未来决策需要的全部信息，转移只从更小且已经算清的状态来。

\textbf{状态复用。} dp[i][j] 表示 word1 前 i 个字符变成 word2 前 j 个字符的最少操作数。按前缀长度从小到大填表。手算时只改被新输入影响的那一小部分，而不是回头重算完整候选。

\textbf{指针和字段。} 状态字段要能说清：i/j 是前缀长度；dp[i-1][j] 对应删除，dp[i][j-1] 对应插入，dp[i-1][j-1] 对应替换或字符相等继承。手算时按这个协议跟踪：空串到 ac 需要插入 2 次；ab 到空串需要删除 2 次。内部格子按末尾字符是否相等决定转移。

\textbf{代码落点。} 关键是第一行 dp[0][j]=j、第一列 dp[i][0]=i。字符相等时不要额外加一。关键代码行必须能逐句翻译成“旧状态怎样变成新状态”，否则就只是模板。

\textbf{正确性。} 任意最优编辑序列的最后一步只有插入、删除、替换或不操作四类，转移枚举了全部最后一步。

\textbf{边界实验。} 先用空串、完全相同、完全不同、操作代价不同做反例检查。实验时覆盖空串、完全相同、完全不同、操作代价不同，先打印完整 DP 表，再比较空间压缩版本。若压缩版不同，优先检查遍历顺序和旧值保存。

\textbf{复杂度变化。} 暴力递归指数级；二维 DP O(mn) 时间和空间，可滚动到 O(min(m, n)) 空间。

\textbf{迁移追问。} 如果题目改成“若只允许插入删除，问题退化到和最长公共子序列相关”，先判断原状态是否仍然覆盖新答案集合，再决定是微调边界、改 value 语义，还是换算法。

\noindent\textbf{LeetCode 85 最大矩形。}

\textbf{题目说明。} LeetCode 85 最大矩形 的题面入口要先还原成具体任务：每一行都可以把连续的一当成柱状图高度。

\textbf{样例入口。} 拿矩阵每一行向上累计高度手算，某行高度数组若为 [2, 1, 2]，就退化成柱状图最大矩形。

\textbf{暴力基线。} 慢版本让每个位置向左或向右线性寻找边界，或者让每个结构反复等待匹配。它能算出答案，但很多尚未完成的候选会被未来同一个事件一起解决。放到本题就是：先按“每一行都可以把连续的一当成柱状图高度”完整试慢版本；样例里已经能看到“规律是逐行更新每列连续 1 的高度，再对每行高度数组跑单调栈；全零行会把高度清零”，所以优化前必须先能说清慢版本枚举了哪些候选。

\textbf{暴力过程。} 拿矩阵每一行向上累计高度手算，某行高度数组若为 [2, 1, 2]，就退化成柱状图最大矩形。

\textbf{重复工作。} 题面模型：每一行都可以把连续的一当成柱状图高度。暴力版的慢点要落到本题：慢版本会围绕“每一行都可以把连续的一当成柱状图高度”反复寻找最近边界或最近未完成对象。样例规律是“规律是逐行更新每列连续 1 的高度，再对每行高度数组跑单调栈；全零行会把高度清零”，说明旧候选一旦遇到当前输入就能被批量结算；继续为每个位置单独向左或向右扫描就是浪费。后面的优化只消掉这类重复工作，不能改变答案集合。

\textbf{优化条件。} 本题的关键观察是：逐行更新每列高度，然后把当前行转成柱状图最大矩形问题。这句话必须能翻译成可维护状态；如果题目条件改变后观察不再成立，就不能继续套原来的写法。

\textbf{相邻状态。} 规律是逐行更新每列连续 1 的高度，再对每行高度数组跑单调栈；全零行会把高度清零。把这个特例继续拆开看：栈里的对象都是答案尚未确定的候选；一旦当前元素触发弹栈，被弹出的候选必须已经同时拿到了左边界和右边界，未来元素不再有资格改变它的答案。

\textbf{状态复用。} 把反复寻找最近边界改成维护未完成候选；新元素只和栈顶比较，连续弹出一批已经被当前元素解决的对象。本题落点是：规律是逐行更新每列连续 1 的高度，再对每行高度数组跑单调栈；全零行会把高度清零。手算时只改被新输入影响的那一小部分，而不是回头重算完整候选。

\textbf{指针和字段。} 状态字段要能说清：栈内对象的业务含义、当前输入、弹出时结算的对象、左边界和右边界；栈中存下标还是值要明确。手算时按这个协议跟踪：拿矩阵每一行向上累计高度手算，某行高度数组若为 [2, 1, 2]，就退化成柱状图最大矩形。然后按协议复查：手算时记录栈内元素的业务含义，不只记录数值。入栈表示答案未定，弹栈表示当前事件已经让它的答案定下来。

\textbf{代码落点。} 弹栈前确认栈非空，弹出后再读取新的左边界；相等元素和哨兵高度要有统一策略。关键代码行必须能逐句翻译成“旧状态怎样变成新状态”，否则就只是模板。

\textbf{正确性。} 证明从弹出时机开始：被弹出的对象在当前时刻答案已经确定，未来元素无法再改变它。

\textbf{边界实验。} 先用全零、全一、单行、单列、宽度和高度结算做反例检查。实验时准备全零、全一、单行、单列、宽度和高度结算，打印每次入栈、弹栈和结算对象。重点观察相等元素、空栈、哨兵和最后一次结算。

\textbf{复杂度变化。} 暴力为每个位置寻找边界常见为平方级；单调栈让每个元素最多入栈一次、出栈一次，因此主体扫描为线性时间。

\textbf{迁移追问。} 如果题目改成“这是二维问题压成一维单调栈的典型做法”，先判断原状态是否仍然覆盖新答案集合，再决定是微调边界、改 value 语义，还是换算法。

\noindent\textbf{LeetCode 97 交错字符串。}

\textbf{题目说明。} LeetCode 97 交错字符串 的题面入口要先还原成具体任务：目标串前缀由两个来源字符串的前缀交错组成。

\textbf{样例入口。} 拿 s1=a、s2=b、s3=ab 手算，最后一个字符 b 可以来自 s2；拿 s1=a、s2=a、s3=aa，两个来源都能匹配，不能贪心固定选一个。

\textbf{暴力基线。} 慢版本先写递归搜索树：节点表示处理位置、剩余资源和当前选择。只要不同路径反复到达同一个节点，就说明需要把这个节点命名成 DP 状态。放到本题就是：先按“目标串前缀由两个来源字符串的前缀交错组成”完整试慢版本；样例里已经能看到“规律是 dp[i][j] 表示两个前缀能否组成目标前 i+j 个字符，转移同时尝试来自 s1 和 s2 的最后一步”，所以优化前必须先能说清慢版本枚举了哪些候选。

\textbf{暴力过程。} 拿 s1=a、s2=b、s3=ab 手算，最后一个字符 b 可以来自 s2；拿 s1=a、s2=a、s3=aa，两个来源都能匹配，不能贪心固定选一个。

\textbf{重复工作。} 题面模型：目标串前缀由两个来源字符串的前缀交错组成。暴力版的慢点要落到本题：慢版本会在“目标串前缀由两个来源字符串的前缀交错组成”的选择树里反复到达相同参数。样例规律是“规律是 dp[i][j] 表示两个前缀能否组成目标前 i+j 个字符，转移同时尝试来自 s1 和 s2 的最后一步”，说明这个参数组合可以命名成状态，算过之后后面直接复用。后面的优化只消掉这类重复工作，不能改变答案集合。

\textbf{优化条件。} 本题的关键观察是：状态 dp[i][j] 表示 s1 前 i 个字符和 s2 前 j 个字符能否组成 s3 前 i 加 j 个字符。这句话必须能翻译成可维护状态；如果题目条件改变后观察不再成立，就不能继续套原来的写法。

\textbf{相邻状态。} 规律是 dp[i][j] 表示两个前缀能否组成目标前 i+j 个字符，转移同时尝试来自 s1 和 s2 的最后一步。把这个特例继续拆开看：先问暴力搜索里哪些不同路径到达同一个后续问题，再给这个后续问题命名为状态。状态必须包含未来决策需要的全部信息，转移只从更小且已经算清的状态来。

\textbf{状态复用。} 把重复递归节点命名为状态；遍历顺序只是在保证依赖状态已经算完。本题落点是：规律是 dp[i][j] 表示两个前缀能否组成目标前 i+j 个字符，转移同时尝试来自 s1 和 s2 的最后一步。手算时只改被新输入影响的那一小部分，而不是回头重算完整候选。

\textbf{指针和字段。} 状态字段要能说清：状态下标、状态值语义、初始化、转移来源、遍历顺序；空间压缩时还要指出读到的是旧值还是新值。手算时按这个协议跟踪：拿 s1=a、s2=b、s3=ab 手算，最后一个字符 b 可以来自 s2；拿 s1=a、s2=a、s3=aa，两个来源都能匹配，不能贪心固定选一个。然后按协议复查：手算时先写一个很小的 DP 表或记忆化搜索记录。每个格子旁边写它来自哪些更小状态。

\textbf{代码落点。} 先写未压缩版本校验转移；压缩前后用小表对拍；不可达哨兵参与加法前要检查溢出。关键代码行必须能逐句翻译成“旧状态怎样变成新状态”，否则就只是模板。

\textbf{正确性。} 证明从最后一步开始：任意最优解的最后一步都在转移枚举中，转移来源已经由更小状态正确计算。

\textbf{边界实验。} 先用总长度不等、第一行第一列初始化、两个来源字符都能匹配时不能贪心只选一个做反例检查。实验时覆盖总长度不等、第一行第一列初始化、两个来源字符都能匹配时不能贪心只选一个，先打印完整 DP 表，再比较空间压缩版本。若压缩版不同，优先检查遍历顺序和旧值保存。

\textbf{复杂度变化。} 暴力递归会重复计算相同子问题，常见指数级；DP 把每个状态算一次，时间是状态数乘单个转移成本，空间是状态表大小或压缩后的大小。

\textbf{迁移追问。} 如果题目改成“若要统计交错方案数，布尔状态要改成计数状态并处理大数或取模”，先判断原状态是否仍然覆盖新答案集合，再决定是微调边界、改 value 语义，还是换算法。

\noindent\textbf{LeetCode 115 不同的子序列。}

\textbf{题目说明。} LeetCode 115 不同的子序列 的题面入口要先还原成具体任务：从源串中删除若干字符后形成目标串，本质是两个前缀之间的计数问题。

\textbf{样例入口。} 拿 s=bab、t=ba 手算。扫描到最后一个 b 时，它不能匹配 t 的 a，只能继承不用它的方案；扫描到 a 时，有两类方案：用这个 a 匹配 t 的末尾，或不用它。

\textbf{暴力基线。} 慢版本先写递归搜索树：节点表示处理位置、剩余资源和当前选择。只要不同路径反复到达同一个节点，就说明需要把这个节点命名成 DP 状态。放到本题就是：先按“从源串中删除若干字符后形成目标串，本质是两个前缀之间的计数问题”完整试慢版本；样例里已经能看到“规律是字符相等时计数相加，不等时只继承不用源字符的方案”，所以优化前必须先能说清慢版本枚举了哪些候选。

\textbf{暴力过程。} 拿 s=bab、t=ba 手算。扫描到最后一个 b 时，它不能匹配 t 的 a，只能继承不用它的方案；扫描到 a 时，有两类方案：用这个 a 匹配 t 的末尾，或不用它。

\textbf{重复工作。} 题面模型：从源串中删除若干字符后形成目标串，本质是两个前缀之间的计数问题。暴力版的慢点要落到本题：慢版本会在“从源串中删除若干字符后形成目标串，本质是两个前缀之间的计数问题”的选择树里反复到达相同参数。样例规律是“规律是字符相等时计数相加，不等时只继承不用源字符的方案”，说明这个参数组合可以命名成状态，算过之后后面直接复用。后面的优化只消掉这类重复工作，不能改变答案集合。

\textbf{优化条件。} 本题的关键观察是：状态 dp[i][j] 表示源串前 i 个字符中有多少个子序列等于目标前 j 个字符。这句话必须能翻译成可维护状态；如果题目条件改变后观察不再成立，就不能继续套原来的写法。

\textbf{相邻状态。} 规律是字符相等时计数相加，不等时只继承不用源字符的方案。把这个特例继续拆开看：先问暴力搜索里哪些不同路径到达同一个后续问题，再给这个后续问题命名为状态。状态必须包含未来决策需要的全部信息，转移只从更小且已经算清的状态来。

\textbf{状态复用。} 把重复递归节点命名为状态；遍历顺序只是在保证依赖状态已经算完。本题落点是：规律是字符相等时计数相加，不等时只继承不用源字符的方案。手算时只改被新输入影响的那一小部分，而不是回头重算完整候选。

\textbf{指针和字段。} 状态字段要能说清：状态下标、状态值语义、初始化、转移来源、遍历顺序；空间压缩时还要指出读到的是旧值还是新值。手算时按这个协议跟踪：拿 s=bab、t=ba 手算。扫描到最后一个 b 时，它不能匹配 t 的 a，只能继承不用它的方案；扫描到 a 时，有两类方案：用这个 a 匹配 t 的末尾，或不用它。然后按协议复查：手算时先写一个很小的 DP 表或记忆化搜索记录。每个格子旁边写它来自哪些更小状态。

\textbf{代码落点。} 先写未压缩版本校验转移；压缩前后用小表对拍；不可达哨兵参与加法前要检查溢出。关键代码行必须能逐句翻译成“旧状态怎样变成新状态”，否则就只是模板。

\textbf{正确性。} 证明从最后一步开始：任意最优解的最后一步都在转移枚举中，转移来源已经由更小状态正确计算。

\textbf{边界实验。} 先用空目标串计数为一、源串为空、重复字符导致计数增长、计数可能超过 int做反例检查。实验时覆盖空目标串计数为一、源串为空、重复字符导致计数增长、计数可能超过 int，先打印完整 DP 表，再比较空间压缩版本。若压缩版不同，优先检查遍历顺序和旧值保存。

\textbf{复杂度变化。} 暴力递归会重复计算相同子问题，常见指数级；DP 把每个状态算一次，时间是状态数乘单个转移成本，空间是状态表大小或压缩后的大小。

\textbf{迁移追问。} 如果题目改成“若只问是否存在目标子序列，可以用双指针，不需要二维计数 DP”，先判断原状态是否仍然覆盖新答案集合，再决定是微调边界、改 value 语义，还是换算法。

\noindent\textbf{LeetCode 126 单词接龙 II。}

\textbf{题目说明。} LeetCode 126 单词接龙 II 的题面入口要先还原成具体任务：不仅要最短转换长度，还要输出所有最短转换路径。

\textbf{样例入口。} 拿 hit 到 cog 的小词表手算，同一层 hot 可以通向 dot 和 lot，后续 dog/log 都可能到 cog。若某个单词在本层第一次见到就立刻删掉，会丢掉另一条同长度父路径。

\textbf{暴力基线。} 慢版本先按题意画出完整状态图，用普通搜索跑通可达性。这个过程用于确认不仅要最短转换长度，还要输出所有最短转换路径的节点和边，之后再讨论访问标记、层次或代价顺序。放到本题就是：先按“不仅要最短转换长度，还要输出所有最短转换路径”完整试慢版本；样例里已经能看到“规律是 BFS 建最短层级，前驱列表保留同层多父节点，最后再回溯所有路径”，所以优化前必须先能说清慢版本枚举了哪些候选。

\textbf{暴力过程。} 拿 hit 到 cog 的小词表手算，同一层 hot 可以通向 dot 和 lot，后续 dog/log 都可能到 cog。若某个单词在本层第一次见到就立刻删掉，会丢掉另一条同长度父路径。

\textbf{重复工作。} 题面模型：不仅要最短转换长度，还要输出所有最短转换路径。暴力版的慢点要落到本题：慢版本会围绕“不仅要最短转换长度，还要输出所有最短转换路径”反复展开同一批状态。样例规律是“规律是 BFS 建最短层级，前驱列表保留同层多父节点，最后再回溯所有路径”，说明必须把节点、边、访问标记和资源维度一次建清；同一状态不应重复入队，不同资源状态也不能误合并。后面的优化只消掉这类重复工作，不能改变答案集合。

\textbf{优化条件。} 本题的关键观察是：BFS 只建立最短层级和前驱列表，同一层的多个前驱都保留，最后再从终点回溯路径。这句话必须能翻译成可维护状态；如果题目条件改变后观察不再成立，就不能继续套原来的写法。

\textbf{相邻状态。} 规律是 BFS 建最短层级，前驱列表保留同层多父节点，最后再回溯所有路径。把这个特例继续拆开看：状态节点不一定等于原始位置，还可能包含钥匙、剩余资源、时间或访问集合。visited 只能合并未来能力完全相同的状态，否则会把合法路径剪掉。

\textbf{状态复用。} 把重复搜索压成访问标记、距离数组或拓扑状态；邻居生成也要从全量扫描变成结构化枚举。本题落点是：规律是 BFS 建最短层级，前驱列表保留同层多父节点，最后再回溯所有路径。手算时只改被新输入影响的那一小部分，而不是回头重算完整候选。

\textbf{指针和字段。} 状态字段要能说清：状态编号、邻接生成方式、访问标记、距离或层数、队列内容；若状态含资源，visited 不能只按位置记录。手算时按这个协议跟踪：拿 hit 到 cog 的小词表手算，同一层 hot 可以通向 dot 和 lot，后续 dog/log 都可能到 cog。若某个单词在本层第一次见到就立刻删掉，会丢掉另一条同长度父路径。然后按协议复查：手算时画队列或栈，状态必须包含题目需要的所有资源。入队时标记还是出队时标记，要和去重语义一致。

\textbf{代码落点。} 入队时标记还是出队时标记必须一致；方向数组集中定义；多源 BFS 要一次性把所有源点放入初始队列。关键代码行必须能逐句翻译成“旧状态怎样变成新状态”，否则就只是模板。

\textbf{正确性。} 证明从可达性开始：搜索覆盖所有合法边能到达的状态，访问标记只删除重复状态而不删除不同未来能力。

\textbf{边界实验。} 先用终点不在词表、同层多父节点、本层访问不能提前删除等长前驱、输出规模可能很大做反例检查。实验时覆盖终点不在词表、同层多父节点、本层访问不能提前删除等长前驱、输出规模可能很大，并打印入队时的完整状态。若同一位置但资源不同的状态被合并，很多图题会出现隐蔽错误。

\textbf{复杂度变化。} 暴力重复从多个起点搜索会反复遍历图；正确建模和访问标记后，BFS 或 DFS 通常按节点数加边数计算，状态图题则按完整状态数计算。

\textbf{迁移追问。} 如果题目改成“若只问最短长度，普通 BFS 或双向 BFS 更轻；若要路径，就必须保存前驱关系”，先判断原状态是否仍然覆盖新答案集合，再决定是微调边界、改 value 语义，还是换算法。

\noindent\textbf{LeetCode 135 分发糖果。}

\textbf{题目说明。} LeetCode 135 分发糖果 的题面入口要先还原成具体任务：相邻评分约束需要同时满足左邻和右邻两个方向。

\textbf{样例入口。} 拿 ratings=[1, 0, 2] 手算，中间孩子评分最低，左右都要比它多糖，结果 [2, 1, 2]。只从左到右会漏掉右侧约束。

\textbf{暴力基线。} 慢版本枚举选择顺序或用 DP 保留多个可能方案，先确认最优值存在。随后才检查局部选择是否能通过交换或领先性固定下来。放到本题就是：先按“相邻评分约束需要同时满足左邻和右邻两个方向”完整试慢版本；样例里已经能看到“规律是左到右满足左邻递增，右到左满足右邻递增，最终取两侧要求最大值”，所以优化前必须先能说清慢版本枚举了哪些候选。

\textbf{暴力过程。} 拿 ratings=[1, 0, 2] 手算，中间孩子评分最低，左右都要比它多糖，结果 [2, 1, 2]。只从左到右会漏掉右侧约束。

\textbf{重复工作。} 题面模型：相邻评分约束需要同时满足左邻和右邻两个方向。暴力版的慢点要落到本题：慢版本会围绕“相邻评分约束需要同时满足左邻和右邻两个方向”枚举选择顺序。样例规律是“规律是左到右满足左邻递增，右到左满足右邻递增，最终取两侧要求最大值”，说明存在一个局部选择或边界状态可以安全保留；不能证明这个选择不变差时，就不能把它叫贪心。后面的优化只消掉这类重复工作，不能改变答案集合。

\textbf{优化条件。} 本题的关键观察是：左到右满足递增关系，右到左满足反向递增关系，最后取两侧要求最大值。这句话必须能翻译成可维护状态；如果题目条件改变后观察不再成立，就不能继续套原来的写法。

\textbf{相邻状态。} 规律是左到右满足左邻递增，右到左满足右邻递增，最终取两侧要求最大值。把这个特例继续拆开看：先构造一个非贪心选择，再比较两者留下的资源、右边界或剩余时间。只有能把非贪心第一步交换成贪心第一步且不变差，局部选择才是规律。

\textbf{状态复用。} 把指数级选择压成一个可证明安全的局部选择；状态通常是当前最远边界、最早结束时间、最小代价或剩余资源。本题落点是：规律是左到右满足左邻递增，右到左满足右邻递增，最终取两侧要求最大值。手算时只改被新输入影响的那一小部分，而不是回头重算完整候选。

\textbf{指针和字段。} 状态字段要能说清：处理顺序、当前已选方案、剩余资源或最远边界、答案计数；若使用排序，要说明排序键；每次局部选择后要能说明当前状态不落后。手算时按这个协议跟踪：拿 ratings=[1, 0, 2] 手算，中间孩子评分最低，左右都要比它多糖，结果 [2, 1, 2]。只从左到右会漏掉右侧约束。然后按协议复查：手算时同时写一个非贪心选择，与贪心选择比较剩余资源。若无法说明贪心后剩余资源不差，就不能直接下结论。

\textbf{代码落点。} 把局部规则写成业务含义；若需要排序，就处理相同关键字；循环中只维护证明需要的状态，不把局部选择和答案更新混在一起。关键代码行必须能逐句翻译成“旧状态怎样变成新状态”，否则就只是模板。

\textbf{正确性。} 证明从交换或领先性开始：任意最优解都能替换成当前贪心选择而不变差，或贪心状态始终不落后。

\textbf{边界实验。} 先用评分相等、严格递增、严格递减、山峰和谷底做反例检查。实验时除了评分相等、严格递增、严格递减、山峰和谷底，还要故意替换成一个看似合理但错误的局部规则，构造反例。能解释错误贪心为何失败，才算理解正确贪心。

\textbf{复杂度变化。} 暴力枚举选择顺序可能指数级；贪心通常先排序，再线性扫描或用堆维护少量候选。

\textbf{迁移追问。} 如果题目改成“这是局部约束两遍扫描，不能只用一个方向贪心”，先判断原状态是否仍然覆盖新答案集合，再决定是微调边界、改 value 语义，还是换算法。

\noindent\textbf{LeetCode 140 单词拆分 II。}

\textbf{题目说明。} LeetCode 140 单词拆分 II 的题面入口要先还原成具体任务：需要输出所有由字典单词拼出的句子，输出规模本身可能指数级。

\textbf{样例入口。} 拿 catsanddog 手算，前缀 cat 和 cats 都能走，但某些后缀继续切会失败。若直接 DFS，每次都会重复探索同一个后缀。

\textbf{暴力基线。} 慢版本就是完整搜索树：每层列出选择列表，路径到达终止条件时检查答案。它先保证覆盖所有可能，再把明显无效或重复的分支剪掉。放到本题就是：先按“需要输出所有由字典单词拼出的句子，输出规模本身可能指数级”完整试慢版本；样例里已经能看到“规律是先用记忆化记录某个起点之后能生成哪些句子，或先用 DP 剪掉不可拆后缀，再枚举路径”，所以优化前必须先能说清慢版本枚举了哪些候选。

\textbf{暴力过程。} 拿 catsanddog 手算，前缀 cat 和 cats 都能走，但某些后缀继续切会失败。若直接 DFS，每次都会重复探索同一个后缀。

\textbf{重复工作。} 题面模型：需要输出所有由字典单词拼出的句子，输出规模本身可能指数级。暴力版的慢点要落到本题：慢版本会围绕“需要输出所有由字典单词拼出的句子，输出规模本身可能指数级”展开完整搜索树。样例规律是“规律是先用记忆化记录某个起点之后能生成哪些句子，或先用 DP 剪掉不可拆后缀，再枚举路径”，说明一层递归必须对应一个明确决策，剪枝只能删除整棵不可能成功或重复的子树。后面的优化只消掉这类重复工作，不能改变答案集合。

\textbf{优化条件。} 本题的关键观察是：先用 DP 或记忆化判断后缀是否可拆，再 DFS 枚举可行前缀并在末尾拼接句子。这句话必须能翻译成可维护状态；如果题目条件改变后观察不再成立，就不能继续套原来的写法。

\textbf{相邻状态。} 规律是先用记忆化记录某个起点之后能生成哪些句子，或先用 DP 剪掉不可拆后缀，再枚举路径。把这个特例继续拆开看：一层递归对应一个明确决策，路径保存已经做出的选择，选择列表保存仍可能扩展的分支。剪枝前必须能说明被剪掉的整棵子树没有合法答案，而不是只因为它看起来不优。

\textbf{状态复用。} 把完整搜索树加上合法性检查、剩余资源剪枝和同层去重；路径状态进入和退出要成对恢复。本题落点是：规律是先用记忆化记录某个起点之后能生成哪些句子，或先用 DP 剪掉不可拆后缀，再枚举路径。手算时只改被新输入影响的那一小部分，而不是回头重算完整候选。

\textbf{指针和字段。} 状态字段要能说清：路径、起点或使用标记、剩余目标、答案集合、剪枝条件；进入递归和退出递归必须成对恢复。手算时按这个协议跟踪：拿 catsanddog 手算，前缀 cat 和 cats 都能走，但某些后缀继续切会失败。若直接 DFS，每次都会重复探索同一个后缀。然后按协议复查：手算时给每层标出选择列表和路径。剪枝发生前要指出被剪分支为什么已经不可能得到合法答案。

\textbf{代码落点。} 剪枝条件写在扩展前；同层去重不要误写成全局去重；保存答案时复制当前路径。关键代码行必须能逐句翻译成“旧状态怎样变成新状态”，否则就只是模板。

\textbf{正确性。} 证明从搜索树覆盖开始：每个合法答案有且只有一条路径，剪枝只删掉不可能成功的分支。

\textbf{边界实验。} 先用无解后缀、重复单词、路径拼接成本、答案数量主导复杂度做反例检查。实验时覆盖无解后缀、重复单词、路径拼接成本、答案数量主导复杂度，统计搜索节点数量和剪枝次数。重复元素题要打印同一层跳过哪些值，确认没有误删合法路径。

\textbf{复杂度变化。} 暴力搜索树的规模由输出或选择空间决定；剪枝不改变最坏指数级本质，但能删除不可能成功或重复的分支，复杂度要说明输出规模。

\textbf{迁移追问。} 如果题目改成“只问能否拆分时用布尔 DP 即可，不要承担枚举所有句子的成本”，先判断原状态是否仍然覆盖新答案集合，再决定是微调边界、改 value 语义，还是换算法。

\noindent\textbf{LeetCode 188 买卖股票 IV。}

\textbf{题目说明。} LeetCode 188 买卖股票 IV 的题面入口要先还原成具体任务：最多 k 次交易让状态必须同时包含交易次数和是否持有股票。

\textbf{样例入口。} 拿 k=2、价格 [3, 2, 6] 手算。某天结束时只说最大利润不够，因为未来能否卖出取决于是否持股，也取决于已经完成几次交易。

\textbf{暴力基线。} 慢版本先写递归搜索树：节点表示处理位置、剩余资源和当前选择。只要不同路径反复到达同一个节点，就说明需要把这个节点命名成 DP 状态。放到本题就是：先按“最多 k 次交易让状态必须同时包含交易次数和是否持有股票”完整试慢版本；样例里已经能看到“规律是状态必须包含交易次数和持有状态；k 很大时才退化为无限次交易”，所以优化前必须先能说清慢版本枚举了哪些候选。

\textbf{暴力过程。} 拿 k=2、价格 [3, 2, 6] 手算。某天结束时只说最大利润不够，因为未来能否卖出取决于是否持股，也取决于已经完成几次交易。

\textbf{重复工作。} 题面模型：最多 k 次交易让状态必须同时包含交易次数和是否持有股票。暴力版的慢点要落到本题：慢版本会在“最多 k 次交易让状态必须同时包含交易次数和是否持有股票”的选择树里反复到达相同参数。样例规律是“规律是状态必须包含交易次数和持有状态；k 很大时才退化为无限次交易”，说明这个参数组合可以命名成状态，算过之后后面直接复用。后面的优化只消掉这类重复工作，不能改变答案集合。

\textbf{优化条件。} 本题的关键观察是：维护每个交易次数下的 buy 和 sell 状态；当 k 足够大时可退化为无限次交易。这句话必须能翻译成可维护状态；如果题目条件改变后观察不再成立，就不能继续套原来的写法。

\textbf{相邻状态。} 规律是状态必须包含交易次数和持有状态；k 很大时才退化为无限次交易。把这个特例继续拆开看：先问暴力搜索里哪些不同路径到达同一个后续问题，再给这个后续问题命名为状态。状态必须包含未来决策需要的全部信息，转移只从更小且已经算清的状态来。

\textbf{状态复用。} 把重复递归节点命名为状态；遍历顺序只是在保证依赖状态已经算完。本题落点是：规律是状态必须包含交易次数和持有状态；k 很大时才退化为无限次交易。手算时只改被新输入影响的那一小部分，而不是回头重算完整候选。

\textbf{指针和字段。} 状态字段要能说清：状态下标、状态值语义、初始化、转移来源、遍历顺序；空间压缩时还要指出读到的是旧值还是新值。手算时按这个协议跟踪：拿 k=2、价格 [3, 2, 6] 手算。某天结束时只说最大利润不够，因为未来能否卖出取决于是否持股，也取决于已经完成几次交易。然后按协议复查：手算时先写一个很小的 DP 表或记忆化搜索记录。每个格子旁边写它来自哪些更小状态。

\textbf{代码落点。} 先写未压缩版本校验转移；压缩前后用小表对拍；不可达哨兵参与加法前要检查溢出。关键代码行必须能逐句翻译成“旧状态怎样变成新状态”，否则就只是模板。

\textbf{正确性。} 证明从最后一步开始：任意最优解的最后一步都在转移枚举中，转移来源已经由更小状态正确计算。

\textbf{边界实验。} 先用k 为零、价格天数很少、k 大于等于 n/2、更新顺序造成同一天状态污染做反例检查。实验时覆盖k 为零、价格天数很少、k 大于等于 n/2、更新顺序造成同一天状态污染，先打印完整 DP 表，再比较空间压缩版本。若压缩版不同，优先检查遍历顺序和旧值保存。

\textbf{复杂度变化。} 暴力递归会重复计算相同子问题，常见指数级；DP 把每个状态算一次，时间是状态数乘单个转移成本，空间是状态表大小或压缩后的大小。

\textbf{迁移追问。} 如果题目改成“手续费、冷冻期和最多两次交易都是同一状态机框架的不同约束”，先判断原状态是否仍然覆盖新答案集合，再决定是微调边界、改 value 语义，还是换算法。

\noindent\textbf{LeetCode 212 单词搜索 II。}

\textbf{题目说明。} LeetCode 212 单词搜索 II 的题面入口要先还原成具体任务：多个单词在同一棋盘上搜索，公共前缀可以共享。

\textbf{样例入口。} 拿 board 上有 oath、pea、eat 手算，单词共享前缀时不应对每个词重复全图搜索；从 o 出发沿 Trie 前缀 o-a-t-h 找到 oath。

\textbf{暴力基线。} 慢版本就是完整搜索树：每层列出选择列表，路径到达终止条件时检查答案。它先保证覆盖所有可能，再把明显无效或重复的分支剪掉。放到本题就是：先按“多个单词在同一棋盘上搜索，公共前缀可以共享”完整试慢版本；样例里已经能看到“规律是 Trie 合并所有单词前缀，DFS 走棋盘同时走 Trie，命中单词后记录并可去重”，所以优化前必须先能说清慢版本枚举了哪些候选。

\textbf{暴力过程。} 拿 board 上有 oath、pea、eat 手算，单词共享前缀时不应对每个词重复全图搜索；从 o 出发沿 Trie 前缀 o-a-t-h 找到 oath。

\textbf{重复工作。} 题面模型：多个单词在同一棋盘上搜索，公共前缀可以共享。暴力版的慢点要落到本题：慢版本会围绕“多个单词在同一棋盘上搜索，公共前缀可以共享”展开完整搜索树。样例规律是“规律是 Trie 合并所有单词前缀，DFS 走棋盘同时走 Trie，命中单词后记录并可去重”，说明一层递归必须对应一个明确决策，剪枝只能删除整棵不可能成功或重复的子树。后面的优化只消掉这类重复工作，不能改变答案集合。

\textbf{优化条件。} 本题的关键观察是：先建 Trie，再从每个格子回溯；路径到达单词结束节点就收集答案。这句话必须能翻译成可维护状态；如果题目条件改变后观察不再成立，就不能继续套原来的写法。

\textbf{相邻状态。} 规律是 Trie 合并所有单词前缀，DFS 走棋盘同时走 Trie，命中单词后记录并可去重。把这个特例继续拆开看：一层递归对应一个明确决策，路径保存已经做出的选择，选择列表保存仍可能扩展的分支。剪枝前必须能说明被剪掉的整棵子树没有合法答案，而不是只因为它看起来不优。

\textbf{状态复用。} 把完整搜索树加上合法性检查、剩余资源剪枝和同层去重；路径状态进入和退出要成对恢复。本题落点是：规律是 Trie 合并所有单词前缀，DFS 走棋盘同时走 Trie，命中单词后记录并可去重。手算时只改被新输入影响的那一小部分，而不是回头重算完整候选。

\textbf{指针和字段。} 状态字段要能说清：路径、起点或使用标记、剩余目标、答案集合、剪枝条件；进入递归和退出递归必须成对恢复。手算时按这个协议跟踪：拿 board 上有 oath、pea、eat 手算，单词共享前缀时不应对每个词重复全图搜索；从 o 出发沿 Trie 前缀 o-a-t-h 找到 oath。然后按协议复查：手算时给每层标出选择列表和路径。剪枝发生前要指出被剪分支为什么已经不可能得到合法答案。

\textbf{代码落点。} 剪枝条件写在扩展前；同层去重不要误写成全局去重；保存答案时复制当前路径。关键代码行必须能逐句翻译成“旧状态怎样变成新状态”，否则就只是模板。

\textbf{正确性。} 证明从搜索树覆盖开始：每个合法答案有且只有一条路径，剪枝只删掉不可能成功的分支。

\textbf{边界实验。} 先用重复单词、同一格不能复用、前缀是完整单词、剪枝删除空子树做反例检查。实验时覆盖重复单词、同一格不能复用、前缀是完整单词、剪枝删除空子树，统计搜索节点数量和剪枝次数。重复元素题要打印同一层跳过哪些值，确认没有误删合法路径。

\textbf{复杂度变化。} 暴力搜索树的规模由输出或选择空间决定；剪枝不改变最坏指数级本质，但能删除不可能成功或重复的分支，复杂度要说明输出规模。

\textbf{迁移追问。} 如果题目改成“这是单词搜索从单目标扩展到多目标后的结构升级”，先判断原状态是否仍然覆盖新答案集合，再决定是微调边界、改 value 语义，还是换算法。

\noindent\textbf{LeetCode 224 基本计算器。}

\textbf{题目说明。} LeetCode 224 基本计算器 的题面入口要先还原成具体任务：括号会开启新的表达式上下文，括号结束时要回到上一层。

\textbf{样例入口。} 拿 1-(2+3) 手算，进入括号前当前结果是 1，符号是 -；括号内算出 5 后要乘以上层符号并合并成 -4。

\textbf{暴力基线。} 慢版本让每个位置向左或向右线性寻找边界，或者让每个结构反复等待匹配。它能算出答案，但很多尚未完成的候选会被未来同一个事件一起解决。放到本题就是：先按“括号会开启新的表达式上下文，括号结束时要回到上一层”完整试慢版本；样例里已经能看到“规律是栈保存进入括号前的结果和符号，一元负号和空格要按字符流处理”，所以优化前必须先能说清慢版本枚举了哪些候选。

\textbf{暴力过程。} 拿 1-(2+3) 手算，进入括号前当前结果是 1，符号是 -；括号内算出 5 后要乘以上层符号并合并成 -4。

\textbf{重复工作。} 题面模型：括号会开启新的表达式上下文，括号结束时要回到上一层。暴力版的慢点要落到本题：慢版本会围绕“括号会开启新的表达式上下文，括号结束时要回到上一层”反复寻找最近边界或最近未完成对象。样例规律是“规律是栈保存进入括号前的结果和符号，一元负号和空格要按字符流处理”，说明旧候选一旦遇到当前输入就能被批量结算；继续为每个位置单独向左或向右扫描就是浪费。后面的优化只消掉这类重复工作，不能改变答案集合。

\textbf{优化条件。} 本题的关键观察是：栈保存进入括号前的结果和符号，当前层算完后乘以上层符号并合并。这句话必须能翻译成可维护状态；如果题目条件改变后观察不再成立，就不能继续套原来的写法。

\textbf{相邻状态。} 规律是栈保存进入括号前的结果和符号，一元负号和空格要按字符流处理。把这个特例继续拆开看：栈里的对象都是答案尚未确定的候选；一旦当前元素触发弹栈，被弹出的候选必须已经同时拿到了左边界和右边界，未来元素不再有资格改变它的答案。

\textbf{状态复用。} 把反复寻找最近边界改成维护未完成候选；新元素只和栈顶比较，连续弹出一批已经被当前元素解决的对象。本题落点是：规律是栈保存进入括号前的结果和符号，一元负号和空格要按字符流处理。手算时只改被新输入影响的那一小部分，而不是回头重算完整候选。

\textbf{指针和字段。} 状态字段要能说清：栈内对象的业务含义、当前输入、弹出时结算的对象、左边界和右边界；栈中存下标还是值要明确。手算时按这个协议跟踪：拿 1-(2+3) 手算，进入括号前当前结果是 1，符号是 -；括号内算出 5 后要乘以上层符号并合并成 -4。然后按协议复查：手算时记录栈内元素的业务含义，不只记录数值。入栈表示答案未定，弹栈表示当前事件已经让它的答案定下来。

\textbf{代码落点。} 弹栈前确认栈非空，弹出后再读取新的左边界；相等元素和哨兵高度要有统一策略。关键代码行必须能逐句翻译成“旧状态怎样变成新状态”，否则就只是模板。

\textbf{正确性。} 证明从弹出时机开始：被弹出的对象在当前时刻答案已经确定，未来元素无法再改变它。

\textbf{边界实验。} 先用多位数、前导空格、一元负号、嵌套括号做反例检查。实验时准备多位数、前导空格、一元负号、嵌套括号，打印每次入栈、弹栈和结算对象。重点观察相等元素、空栈、哨兵和最后一次结算。

\textbf{复杂度变化。} 暴力为每个位置寻找边界常见为平方级；单调栈让每个元素最多入栈一次、出栈一次，因此主体扫描为线性时间。

\textbf{迁移追问。} 如果题目改成“基本计算器 II 没有括号但有乘除优先级，需要不同的栈语义”，先判断原状态是否仍然覆盖新答案集合，再决定是微调边界、改 value 语义，还是换算法。

\noindent\textbf{LeetCode 227 基本计算器 II。}

\textbf{题目说明。} LeetCode 227 基本计算器 II 的题面入口要先还原成具体任务：乘除优先级高于加减，可以在读到下一个运算符时结算上一项。

\textbf{样例入口。} 拿 3+2*2 手算，加号前的 3 可以先入栈，乘号遇到 2 时要立刻把栈顶 2 乘以下一个数 2，而不能等到最后按顺序相加。

\textbf{暴力基线。} 慢版本让每个位置向左或向右线性寻找边界，或者让每个结构反复等待匹配。它能算出答案，但很多尚未完成的候选会被未来同一个事件一起解决。放到本题就是：先按“乘除优先级高于加减，可以在读到下一个运算符时结算上一项”完整试慢版本；样例里已经能看到“规律是栈保存已确定符号的项，乘除立即修改栈顶，加减压入正负项”，所以优化前必须先能说清慢版本枚举了哪些候选。

\textbf{暴力过程。} 拿 3+2*2 手算，加号前的 3 可以先入栈，乘号遇到 2 时要立刻把栈顶 2 乘以下一个数 2，而不能等到最后按顺序相加。

\textbf{重复工作。} 题面模型：乘除优先级高于加减，可以在读到下一个运算符时结算上一项。暴力版的慢点要落到本题：慢版本会围绕“乘除优先级高于加减，可以在读到下一个运算符时结算上一项”反复寻找最近边界或最近未完成对象。样例规律是“规律是栈保存已确定符号的项，乘除立即修改栈顶，加减压入正负项”，说明旧候选一旦遇到当前输入就能被批量结算；继续为每个位置单独向左或向右扫描就是浪费。后面的优化只消掉这类重复工作，不能改变答案集合。

\textbf{优化条件。} 本题的关键观察是：栈保存已经确定符号的项，乘除直接修改栈顶，加减压入正负项。这句话必须能翻译成可维护状态；如果题目条件改变后观察不再成立，就不能继续套原来的写法。

\textbf{相邻状态。} 规律是栈保存已确定符号的项，乘除立即修改栈顶，加减压入正负项。把这个特例继续拆开看：栈里的对象都是答案尚未确定的候选；一旦当前元素触发弹栈，被弹出的候选必须已经同时拿到了左边界和右边界，未来元素不再有资格改变它的答案。

\textbf{状态复用。} 把反复寻找最近边界改成维护未完成候选；新元素只和栈顶比较，连续弹出一批已经被当前元素解决的对象。本题落点是：规律是栈保存已确定符号的项，乘除立即修改栈顶，加减压入正负项。手算时只改被新输入影响的那一小部分，而不是回头重算完整候选。

\textbf{指针和字段。} 状态字段要能说清：栈内对象的业务含义、当前输入、弹出时结算的对象、左边界和右边界；栈中存下标还是值要明确。手算时按这个协议跟踪：拿 3+2*2 手算，加号前的 3 可以先入栈，乘号遇到 2 时要立刻把栈顶 2 乘以下一个数 2，而不能等到最后按顺序相加。然后按协议复查：手算时记录栈内元素的业务含义，不只记录数值。入栈表示答案未定，弹栈表示当前事件已经让它的答案定下来。

\textbf{代码落点。} 弹栈前确认栈非空，弹出后再读取新的左边界；相等元素和哨兵高度要有统一策略。关键代码行必须能逐句翻译成“旧状态怎样变成新状态”，否则就只是模板。

\textbf{正确性。} 证明从弹出时机开始：被弹出的对象在当前时刻答案已经确定，未来元素无法再改变它。

\textbf{边界实验。} 先用多位数、空格、除法向零截断、最后一个数字需要结算做反例检查。实验时准备多位数、空格、除法向零截断、最后一个数字需要结算，打印每次入栈、弹栈和结算对象。重点观察相等元素、空栈、哨兵和最后一次结算。

\textbf{复杂度变化。} 暴力为每个位置寻找边界常见为平方级；单调栈让每个元素最多入栈一次、出栈一次，因此主体扫描为线性时间。

\textbf{迁移追问。} 如果题目改成“若加入括号，需要再引入上下文栈或递归下降解析”，先判断原状态是否仍然覆盖新答案集合，再决定是微调边界、改 value 语义，还是换算法。

\noindent\textbf{LeetCode 295 数据流的中位数。}

\textbf{题目说明。} LeetCode 295 数据流的中位数 的题面入口要先还原成具体任务：数据流需要动态维护较小一半和较大一半。

\textbf{样例入口。} 拿数据流 1、2、3 手算，插入 1 后中位数 1；插入 2 后两半是 [1] 和 [2]，中位数 1.5；插入 3 后右半多一个，要平衡成左半 [2, 1]、右半 [3]。

\textbf{暴力基线。} 暴力每插入一个数就重新排序所有已到达元素，然后按长度奇偶取中位数。

\textbf{暴力过程。} 数据流 1、2、3 中，插入 1 中位数是 1；插入 2 后中位数是 (1+2)/2；插入 3 后中位数是 2。

\textbf{重复工作。} 题面模型：数据流需要动态维护较小一半和较大一半。暴力版的慢点要落到本题：浪费在每次都维护完整排序。中位数只需要较小一半的最大值和较大一半的最小值。后面的优化只消掉这类重复工作，不能改变答案集合。

\textbf{优化条件。} 本题的关键观察是：大顶堆保存较小一半，小顶堆保存较大一半，大小差不超过一。这句话必须能翻译成可维护状态；如果题目条件改变后观察不再成立，就不能继续套原来的写法。

\textbf{相邻状态。} 规律是大顶堆保存较小一半，小顶堆保存较大一半，大小差最多 1。把这个特例继续拆开看：堆顶必须正好代表下一步最需要处理的候选；若堆里可能有过期对象，读取堆顶前要先清理。堆只解决动态极值，不自动证明被弹出或被丢弃的元素无用。

\textbf{状态复用。} 大顶堆 small 保存较小一半，小顶堆 large 保存较大一半。保持 small.size() 与 large.size() 差不超过 1，且 small.top()<=large.top()。手算时只改被新输入影响的那一小部分，而不是回头重算完整候选。

\textbf{指针和字段。} 状态字段要能说清：small.top 是左半最大值，large.top 是右半最小值；奇数时较大堆顶是中位数，偶数时两个堆顶平均。手算时按这个协议跟踪：插入 1 放 small；插入 2 后移动平衡成 small=[1]、large=[2]；插入 3 后 large 多，再移动 2 到 small，中位数 2。

\textbf{代码落点。} 关键是先按顺序不变量放入堆，再重平衡大小。求平均时用 double，并防止两个 int 相加溢出。关键代码行必须能逐句翻译成“旧状态怎样变成新状态”，否则就只是模板。

\textbf{正确性。} 两个堆按大小分割数据流，中位数只可能位于左半最大或右半最小；大小不变量保证取堆顶即可。

\textbf{边界实验。} 先用偶数个元素、负数、重复值、平均值溢出做反例检查。实验时覆盖偶数个元素、负数、重复值、平均值溢出，记录堆顶是否过期、比较器是否让正确候选在堆顶、K 或窗口大小为边界值时是否稳定。

\textbf{复杂度变化。} 每次重新排序 O(n log n)；双堆 addNum O(log n)，findMedian O(1)，空间 O(n)。

\textbf{迁移追问。} 如果题目改成“若还要删除任意元素，需要延迟删除哈希表”，先判断原状态是否仍然覆盖新答案集合，再决定是微调边界、改 value 语义，还是换算法。

\noindent\textbf{LeetCode 312 戳气球。}

\textbf{题目说明。} LeetCode 312 戳气球 的题面入口要先还原成具体任务：先戳哪个会让邻居持续变化，改成最后戳某个气球后区间边界才稳定。

\textbf{样例入口。} 拿 [3, 1, 5] 手算，若先戳 1，左右邻居会变；若改问某个区间里最后戳谁，最后一刻左右边界固定，收益可拆成左右两个子区间。

\textbf{暴力基线。} 慢版本先写递归搜索树：节点表示处理位置、剩余资源和当前选择。只要不同路径反复到达同一个节点，就说明需要把这个节点命名成 DP 状态。放到本题就是：先按“先戳哪个会让邻居持续变化，改成最后戳某个气球后区间边界才稳定”完整试慢版本；样例里已经能看到“规律是区间 DP 枚举最后戳的 mid，而不是第一个戳的气球”，所以优化前必须先能说清慢版本枚举了哪些候选。

\textbf{暴力过程。} 拿 [3, 1, 5] 手算，若先戳 1，左右邻居会变；若改问某个区间里最后戳谁，最后一刻左右边界固定，收益可拆成左右两个子区间。

\textbf{重复工作。} 题面模型：先戳哪个会让邻居持续变化，改成最后戳某个气球后区间边界才稳定。暴力版的慢点要落到本题：慢版本会在“先戳哪个会让邻居持续变化，改成最后戳某个气球后区间边界才稳定”的选择树里反复到达相同参数。样例规律是“规律是区间 DP 枚举最后戳的 mid，而不是第一个戳的气球”，说明这个参数组合可以命名成状态，算过之后后面直接复用。后面的优化只消掉这类重复工作，不能改变答案集合。

\textbf{优化条件。} 本题的关键观察是：区间 DP 定义开区间内全部戳完的最大收益，枚举最后一个被戳的气球作为切分点。这句话必须能翻译成可维护状态；如果题目条件改变后观察不再成立，就不能继续套原来的写法。

\textbf{相邻状态。} 规律是区间 DP 枚举最后戳的 mid，而不是第一个戳的气球。把这个特例继续拆开看：先问暴力搜索里哪些不同路径到达同一个后续问题，再给这个后续问题命名为状态。状态必须包含未来决策需要的全部信息，转移只从更小且已经算清的状态来。

\textbf{状态复用。} 把重复递归节点命名为状态；遍历顺序只是在保证依赖状态已经算完。本题落点是：规律是区间 DP 枚举最后戳的 mid，而不是第一个戳的气球。手算时只改被新输入影响的那一小部分，而不是回头重算完整候选。

\textbf{指针和字段。} 状态字段要能说清：状态下标、状态值语义、初始化、转移来源、遍历顺序；空间压缩时还要指出读到的是旧值还是新值。手算时按这个协议跟踪：拿 [3, 1, 5] 手算，若先戳 1，左右邻居会变；若改问某个区间里最后戳谁，最后一刻左右边界固定，收益可拆成左右两个子区间。然后按协议复查：手算时先写一个很小的 DP 表或记忆化搜索记录。每个格子旁边写它来自哪些更小状态。

\textbf{代码落点。} 先写未压缩版本校验转移；压缩前后用小表对拍；不可达哨兵参与加法前要检查溢出。关键代码行必须能逐句翻译成“旧状态怎样变成新状态”，否则就只是模板。

\textbf{正确性。} 证明从最后一步开始：任意最优解的最后一步都在转移枚举中，转移来源已经由更小状态正确计算。

\textbf{边界实验。} 先用补一的虚拟边界、区间长度遍历顺序、空区间收益为零、乘积可能较大做反例检查。实验时覆盖补一的虚拟边界、区间长度遍历顺序、空区间收益为零、乘积可能较大，先打印完整 DP 表，再比较空间压缩版本。若压缩版不同，优先检查遍历顺序和旧值保存。

\textbf{复杂度变化。} 暴力递归会重复计算相同子问题，常见指数级；DP 把每个状态算一次，时间是状态数乘单个转移成本，空间是状态表大小或压缩后的大小。

\textbf{迁移追问。} 如果题目改成“很多区间 DP 都要换成最后一步思考，避免中间操作改变边界”，先判断原状态是否仍然覆盖新答案集合，再决定是微调边界、改 value 语义，还是换算法。

\noindent\textbf{LeetCode 329 矩阵中的最长递增路径。}

\textbf{题目说明。} LeetCode 329 矩阵中的最长递增路径 的题面入口要先还原成具体任务：每个格子的答案是从它出发能走出的最长严格递增路径。

\textbf{样例入口。} 拿矩阵中 1->2->3 的小路径手算，从 1 出发会用到从 2 出发的答案；如果另一个格子也走到 2，后缀路径被重复计算。严格递增保证不会成环。

\textbf{暴力基线。} 慢版本先写递归搜索树：节点表示处理位置、剩余资源和当前选择。只要不同路径反复到达同一个节点，就说明需要把这个节点命名成 DP 状态。放到本题就是：先按“每个格子的答案是从它出发能走出的最长严格递增路径”完整试慢版本；样例里已经能看到“规律是每个格子记忆化保存从这里出发的最长路径”，所以优化前必须先能说清慢版本枚举了哪些候选。

\textbf{暴力过程。} 拿矩阵中 1->2->3 的小路径手算，从 1 出发会用到从 2 出发的答案；如果另一个格子也走到 2，后缀路径被重复计算。严格递增保证不会成环。

\textbf{重复工作。} 题面模型：每个格子的答案是从它出发能走出的最长严格递增路径。暴力版的慢点要落到本题：慢版本会在“每个格子的答案是从它出发能走出的最长严格递增路径”的选择树里反复到达相同参数。样例规律是“规律是每个格子记忆化保存从这里出发的最长路径”，说明这个参数组合可以命名成状态，算过之后后面直接复用。后面的优化只消掉这类重复工作，不能改变答案集合。

\textbf{优化条件。} 本题的关键观察是：把严格递增关系看成无环状态图，用记忆化 DFS 或按值排序的 DAG DP 复用后续格子结果。这句话必须能翻译成可维护状态；如果题目条件改变后观察不再成立，就不能继续套原来的写法。

\textbf{相邻状态。} 规律是每个格子记忆化保存从这里出发的最长路径。把这个特例继续拆开看：先问暴力搜索里哪些不同路径到达同一个后续问题，再给这个后续问题命名为状态。状态必须包含未来决策需要的全部信息，转移只从更小且已经算清的状态来。

\textbf{状态复用。} 把重复递归节点命名为状态；遍历顺序只是在保证依赖状态已经算完。本题落点是：规律是每个格子记忆化保存从这里出发的最长路径。手算时只改被新输入影响的那一小部分，而不是回头重算完整候选。

\textbf{指针和字段。} 状态字段要能说清：状态下标、状态值语义、初始化、转移来源、遍历顺序；空间压缩时还要指出读到的是旧值还是新值。手算时按这个协议跟踪：拿矩阵中 1->2->3 的小路径手算，从 1 出发会用到从 2 出发的答案；如果另一个格子也走到 2，后缀路径被重复计算。严格递增保证不会成环。然后按协议复查：手算时先写一个很小的 DP 表或记忆化搜索记录。每个格子旁边写它来自哪些更小状态。

\textbf{代码落点。} 先写未压缩版本校验转移；压缩前后用小表对拍；不可达哨兵参与加法前要检查溢出。关键代码行必须能逐句翻译成“旧状态怎样变成新状态”，否则就只是模板。

\textbf{正确性。} 证明从最后一步开始：任意最优解的最后一步都在转移枚举中，转移来源已经由更小状态正确计算。

\textbf{边界实验。} 先用平台相等不能走、四方向越界、递归深度、允许不下降时会产生环做反例检查。实验时覆盖平台相等不能走、四方向越界、递归深度、允许不下降时会产生环，先打印完整 DP 表，再比较空间压缩版本。若压缩版不同，优先检查遍历顺序和旧值保存。

\textbf{复杂度变化。} 暴力递归会重复计算相同子问题，常见指数级；DP 把每个状态算一次，时间是状态数乘单个转移成本，空间是状态表大小或压缩后的大小。

\textbf{迁移追问。} 如果题目改成“若改成求最短路径，严格递增的 DP 模型不再直接适用”，先判断原状态是否仍然覆盖新答案集合，再决定是微调边界、改 value 语义，还是换算法。

\noindent\textbf{LeetCode 394 字符串解码。}

\textbf{题目说明。} LeetCode 394 字符串解码 的题面入口要先还原成具体任务：嵌套结构需要保存上一层字符串和重复次数。

\textbf{样例入口。} 拿 3[a2[c]] 手算，遇到左括号前保存当前字符串和重复次数；遇到右括号时弹出上一层，把当前 cc 重复 2 次拼成 acc，再被外层重复 3 次。

\textbf{暴力基线。} 慢版本让每个位置向左或向右线性寻找边界，或者让每个结构反复等待匹配。它能算出答案，但很多尚未完成的候选会被未来同一个事件一起解决。放到本题就是：先按“嵌套结构需要保存上一层字符串和重复次数”完整试慢版本；样例里已经能看到“规律是栈保存进入括号前的上下文，右括号触发一次完整展开”，所以优化前必须先能说清慢版本枚举了哪些候选。

\textbf{暴力过程。} 拿 3[a2[c]] 手算，遇到左括号前保存当前字符串和重复次数；遇到右括号时弹出上一层，把当前 cc 重复 2 次拼成 acc，再被外层重复 3 次。

\textbf{重复工作。} 题面模型：嵌套结构需要保存上一层字符串和重复次数。暴力版的慢点要落到本题：慢版本会围绕“嵌套结构需要保存上一层字符串和重复次数”反复寻找最近边界或最近未完成对象。样例规律是“规律是栈保存进入括号前的上下文，右括号触发一次完整展开”，说明旧候选一旦遇到当前输入就能被批量结算；继续为每个位置单独向左或向右扫描就是浪费。后面的优化只消掉这类重复工作，不能改变答案集合。

\textbf{优化条件。} 本题的关键观察是：遇到左括号入栈当前状态，右括号弹出并拼接展开。这句话必须能翻译成可维护状态；如果题目条件改变后观察不再成立，就不能继续套原来的写法。

\textbf{相邻状态。} 规律是栈保存进入括号前的上下文，右括号触发一次完整展开。把这个特例继续拆开看：栈里的对象都是答案尚未确定的候选；一旦当前元素触发弹栈，被弹出的候选必须已经同时拿到了左边界和右边界，未来元素不再有资格改变它的答案。

\textbf{状态复用。} 把反复寻找最近边界改成维护未完成候选；新元素只和栈顶比较，连续弹出一批已经被当前元素解决的对象。本题落点是：规律是栈保存进入括号前的上下文，右括号触发一次完整展开。手算时只改被新输入影响的那一小部分，而不是回头重算完整候选。

\textbf{指针和字段。} 状态字段要能说清：栈内对象的业务含义、当前输入、弹出时结算的对象、左边界和右边界；栈中存下标还是值要明确。手算时按这个协议跟踪：拿 3[a2[c]] 手算，遇到左括号前保存当前字符串和重复次数；遇到右括号时弹出上一层，把当前 cc 重复 2 次拼成 acc，再被外层重复 3 次。然后按协议复查：手算时记录栈内元素的业务含义，不只记录数值。入栈表示答案未定，弹栈表示当前事件已经让它的答案定下来。

\textbf{代码落点。} 弹栈前确认栈非空，弹出后再读取新的左边界；相等元素和哨兵高度要有统一策略。关键代码行必须能逐句翻译成“旧状态怎样变成新状态”，否则就只是模板。

\textbf{正确性。} 证明从弹出时机开始：被弹出的对象在当前时刻答案已经确定，未来元素无法再改变它。

\textbf{边界实验。} 先用多位数字、嵌套、空片段、数字后必须跟括号做反例检查。实验时准备多位数字、嵌套、空片段、数字后必须跟括号，打印每次入栈、弹栈和结算对象。重点观察相等元素、空栈、哨兵和最后一次结算。

\textbf{复杂度变化。} 暴力为每个位置寻找边界常见为平方级；单调栈让每个元素最多入栈一次、出栈一次，因此主体扫描为线性时间。

\textbf{迁移追问。} 如果题目改成“递归下降解析更通用，但迭代栈更符合工程栈控制”，先判断原状态是否仍然覆盖新答案集合，再决定是微调边界、改 value 语义，还是换算法。

\noindent\textbf{LeetCode 460 LFU 缓存。}

\textbf{题目说明。} LeetCode 460 LFU 缓存 的题面入口要先还原成具体任务：缓存淘汰同时按访问频次和同频最近性排序，需要维护两个维度的索引。

\textbf{样例入口。} 拿容量 2，put(1)、put(2)、get(1)、put(3) 手算。此时 1 的频次高于 2，应淘汰 2；若两个 key 频次相同，还要淘汰同频里最旧的。

\textbf{暴力基线。} 慢版本可以先把节点顺序抄到数组里，按数组推导目标顺序。回到链表后，难点不再是值的顺序，而是节点身份和 next 指针不能丢。放到本题就是：先按“缓存淘汰同时按访问频次和同频最近性排序，需要维护两个维度的索引”完整试慢版本；样例里已经能看到“规律是 key 哈希定位节点，频次哈希定位同频链表，minFreq 指向当前最低频次”，所以优化前必须先能说清慢版本枚举了哪些候选。

\textbf{暴力过程。} 拿容量 2，put(1)、put(2)、get(1)、put(3) 手算。此时 1 的频次高于 2，应淘汰 2；若两个 key 频次相同，还要淘汰同频里最旧的。

\textbf{重复工作。} 题面模型：缓存淘汰同时按访问频次和同频最近性排序，需要维护两个维度的索引。暴力版的慢点要落到本题：慢版本会围绕“缓存淘汰同时按访问频次和同频最近性排序，需要维护两个维度的索引”借助数组或多次扫描确认目标顺序。样例规律是“规律是 key 哈希定位节点，频次哈希定位同频链表，minFreq 指向当前最低频次”，说明优化点不在节点值，而在少量指针角色能否保持已处理链、未处理链和接回位置都不丢。后面的优化只消掉这类重复工作，不能改变答案集合。

\textbf{优化条件。} 本题的关键观察是：哈希表按 key 定位节点，频次到双向链表保存同频节点顺序，并维护当前最小频次。这句话必须能翻译成可维护状态；如果题目条件改变后观察不再成立，就不能继续套原来的写法。

\textbf{相邻状态。} 规律是 key 哈希定位节点，频次哈希定位同频链表，minFreq 指向当前最低频次。把这个特例继续拆开看：每个指针都要有角色，上一段尾、当前段头、当前节点和后继入口不能混。只要一次改 next 之前没有保存后继，就可能丢掉整段未处理链。

\textbf{状态复用。} 把数组辅助结构换成几个稳定指针角色；重连顺序必须保证已处理链合法、未处理链仍可达。本题落点是：规律是 key 哈希定位节点，频次哈希定位同频链表，minFreq 指向当前最低频次。手算时只改被新输入影响的那一小部分，而不是回头重算完整候选。

\textbf{指针和字段。} 状态字段要能说清：虚拟头、上一段尾、当前段头、当前节点、后继节点；任何改 next 的操作之前都要保存后续入口。手算时按这个协议跟踪：拿容量 2，put(1)、put(2)、get(1)、put(3) 手算。此时 1 的频次高于 2，应淘汰 2；若两个 key 频次相同，还要淘汰同频里最旧的。然后按协议复查：手算时画出 prev、cur、next 和下一段头节点。每次改 next 前先保存后继，这是链表推导的安全线。

\textbf{代码落点。} 所有重连步骤按保存后继、断开或反转、接回前缀、接回后缀的顺序写；最后检查是否形成环或丢节点。关键代码行必须能逐句翻译成“旧状态怎样变成新状态”，否则就只是模板。

\textbf{正确性。} 证明从链结构开始：已处理链连通、未处理链可达、二者连接点明确，节点没有丢失也没有成环。

\textbf{边界实验。} 先用容量为零、更新已有 key、频次链表变空时最小频次更新、同频淘汰最旧节点做反例检查。实验时覆盖容量为零、更新已有 key、频次链表变空时最小频次更新、同频淘汰最旧节点，每步画出前驱、当前、后继和下一段头。链表题不要只看输出值，还要检查节点是否丢失或成环。

\textbf{复杂度变化。} 暴力可能多次扫描链表或拷贝到数组；优化后每个节点通常只被访问、重连常数次，目标是线性时间和常数额外空间。

\textbf{迁移追问。} 如果题目改成“LRU 只有时间顺序一个维度，LFU 多了频次维度，结构维护明显更复杂”，先判断原状态是否仍然覆盖新答案集合，再决定是微调边界、改 value 语义，还是换算法。

\noindent\textbf{LeetCode 560 和为 K 的子数组。}

\textbf{题目说明。} LeetCode 560 和为 K 的子数组给定整数数组 nums 和整数 k，统计和等于 k 的连续子数组个数；数组中可以有负数。

\textbf{样例入口。} 样例 nums=[1, 1, 1]，k=2，输出 2；两个合法连续子数组分别是前两个 1 和后两个 1。

\textbf{暴力基线。} 暴力固定 left，向右累计 sum；每当 sum==k 就计数。

\textbf{暴力过程。} 以 nums=[1, 1, 1]、k=2 为例，从 left=0 可找到 [1, 1]，从 left=1 又找到后两个 1。

\textbf{重复工作。} 题面模型：当前前缀和减去目标值，就是需要匹配的历史前缀。暴力版的慢点要落到本题：浪费在每个右端都回头枚举所有 left。区间和等于 k 可以改写成 prefix[right]-prefix[left]=k。后面的优化只消掉这类重复工作，不能改变答案集合。

\textbf{优化条件。} 本题的关键观察是：哈希表保存前缀和出现次数，先查询再更新。这句话必须能翻译成可维护状态；如果题目条件改变后观察不再成立，就不能继续套原来的写法。

\textbf{相邻状态。} 规律是哈希表保存前缀和频次，且初始要放 0:1。把这个特例继续拆开看：每个合法答案都要还原成当前状态和某个历史状态的配对；哈希表的 key 负责定位历史状态，value 负责表达次数、最早位置或存在性。先查还是先写，必须由是否允许当前前缀和自己配对来决定。

\textbf{状态复用。} 扫描前缀和 prefix。哈希表 count[old\_prefix] 保存历史前缀出现次数；当前位置需要查 prefix-k 出现过几次。手算时只改被新输入影响的那一小部分，而不是回头重算完整候选。

\textbf{指针和字段。} 状态字段要能说清：prefix 是当前前缀和，need=prefix-k 是所需历史前缀，count 保存历史前缀频次，answer 累加匹配次数。手算时按这个协议跟踪：[1, 1, 1] 中扫描到第二个元素时 prefix=2，需要历史 0，count[0]=1，所以找到第一个长度 2 子数组。

\textbf{代码落点。} 关键是先 answer += count[prefix-k]，再 count[prefix]++。初始化 count[0]=1，表示从下标 0 开始的区间。关键代码行必须能逐句翻译成“旧状态怎样变成新状态”，否则就只是模板。

\textbf{正确性。} 任意和为 k 的子数组都对应一对前缀差 k；扫描到右端时，所有合法左端前缀都已在表里。

\textbf{边界实验。} 先用负数、目标为零、从下标零开始的区间、重复前缀做反例检查。实验时把负数、目标为零、从下标零开始的区间、重复前缀加入对拍集合，用平方暴力和优化版比较。失败时先看初始历史状态、查询更新顺序和哈希表 value 类型。

\textbf{复杂度变化。} 暴力 O(\ensuremath{n^{2}})，前缀哈希 O(n) 均摊时间、O(n) 空间。

\textbf{迁移追问。} 如果题目改成“若要求最长区间，哈希表应保存最早位置而不是次数”，先判断原状态是否仍然覆盖新答案集合，再决定是微调边界、改 value 语义，还是换算法。

\noindent\textbf{LeetCode 632 最小区间覆盖 K 个列表。}

\textbf{题目说明。} LeetCode 632 最小区间覆盖 K 个列表 的题面入口要先还原成具体任务：每个列表至少取一个元素时，当前区间由最小当前值和最大当前值决定。

\textbf{样例入口。} 拿三组 [4, 10]、[0, 9]、[5, 18] 手算，当前头是 4、0、5，区间 [0, 5] 覆盖三组；弹出最小头 0 后推进该组。

\textbf{暴力基线。} 慢版本每轮扫描全部候选来找最大或最小，再更新候选集合。它的答案选择过程清楚，但动态集合每变化一次就重新找极值，代价太高。放到本题就是：先按“每个列表至少取一个元素时，当前区间由最小当前值和最大当前值决定”完整试慢版本；样例里已经能看到“规律是堆保存每组当前元素，同时维护当前最大值，最小头决定下一步推进哪一组”，所以优化前必须先能说清慢版本枚举了哪些候选。

\textbf{暴力过程。} 拿三组 [4, 10]、[0, 9]、[5, 18] 手算，当前头是 4、0、5，区间 [0, 5] 覆盖三组；弹出最小头 0 后推进该组。

\textbf{重复工作。} 题面模型：每个列表至少取一个元素时，当前区间由最小当前值和最大当前值决定。暴力版的慢点要落到本题：慢版本会围绕“每个列表至少取一个元素时，当前区间由最小当前值和最大当前值决定”反复扫描动态候选集找极值。样例规律是“规律是堆保存每组当前元素，同时维护当前最大值，最小头决定下一步推进哪一组”，说明每轮真正需要的是堆顶边界，而不是把候选全集重新排序。后面的优化只消掉这类重复工作，不能改变答案集合。

\textbf{优化条件。} 本题的关键观察是：小顶堆保存各列表当前元素，同时维护当前最大值；弹出最小所在列表并推进。这句话必须能翻译成可维护状态；如果题目条件改变后观察不再成立，就不能继续套原来的写法。

\textbf{相邻状态。} 规律是堆保存每组当前元素，同时维护当前最大值，最小头决定下一步推进哪一组。把这个特例继续拆开看：堆顶必须正好代表下一步最需要处理的候选；若堆里可能有过期对象，读取堆顶前要先清理。堆只解决动态极值，不自动证明被弹出或被丢弃的元素无用。

\textbf{状态复用。} 把每轮全量扫描改成堆顶查询；插入、弹出和过期清理共同维护动态候选集。本题落点是：规律是堆保存每组当前元素，同时维护当前最大值，最小头决定下一步推进哪一组。手算时只改被新输入影响的那一小部分，而不是回头重算完整候选。

\textbf{指针和字段。} 状态字段要能说清：堆元素字段、堆顶比较规则、候选是否过期、答案读取时机；如果需要懒删除，还要保存删除计数。手算时按这个协议跟踪：拿三组 [4, 10]、[0, 9]、[5, 18] 手算，当前头是 4、0、5，区间 [0, 5] 覆盖三组；弹出最小头 0 后推进该组。然后按协议复查：手算时记录堆顶、候选集合和是否有过期元素。每次使用堆顶前都要确认它仍然属于当前问题。

\textbf{代码落点。} 比较器方向先用小样例验证；每次使用堆顶前清理过期项；堆中保存下标时要能回查原数据。关键代码行必须能逐句翻译成“旧状态怎样变成新状态”，否则就只是模板。

\textbf{正确性。} 证明从候选覆盖开始：所有可能成为下一步答案的对象都在堆中，堆顶过期时不会被使用。

\textbf{边界实验。} 先用某个列表为空、重复值、区间长度相同的 tie-breaker、列表耗尽做反例检查。实验时覆盖某个列表为空、重复值、区间长度相同的 tie-breaker、列表耗尽，记录堆顶是否过期、比较器是否让正确候选在堆顶、K 或窗口大小为边界值时是否稳定。

\textbf{复杂度变化。} 暴力每轮扫描或排序动态候选，会把线性扫描或排序反复发生；堆把取极值、插入和删除边界控制在对数级。

\textbf{迁移追问。} 如果题目改成“这是多路归并和覆盖窗口的结合”，先判断原状态是否仍然覆盖新答案集合，再决定是微调边界、改 value 语义，还是换算法。

\noindent\textbf{LeetCode 678 有效的括号字符串。}

\textbf{题目说明。} LeetCode 678 有效的括号字符串 的题面入口要先还原成具体任务：星号可以作为左括号、右括号或空，使未匹配左括号数量变成一个可能区间。

\textbf{样例入口。} 拿 (*) 手算，星号可以当空或左括号；拿 (*))，星号必须当左括号才能匹配后面的右括号。与其枚举星号三种选择，不如维护可能未匹配左括号数的区间 [low, high]。

\textbf{暴力基线。} 慢版本枚举选择顺序或用 DP 保留多个可能方案，先确认最优值存在。随后才检查局部选择是否能通过交换或领先性固定下来。放到本题就是：先按“星号可以作为左括号、右括号或空，使未匹配左括号数量变成一个可能区间”完整试慢版本；样例里已经能看到“规律是左括号让区间加一，右括号让区间减一，星号让下界减一上界加一，且下界不能低于零”，所以优化前必须先能说清慢版本枚举了哪些候选。

\textbf{暴力过程。} 拿 (*) 手算，星号可以当空或左括号；拿 (*))，星号必须当左括号才能匹配后面的右括号。与其枚举星号三种选择，不如维护可能未匹配左括号数的区间 [low, high]。

\textbf{重复工作。} 题面模型：星号可以作为左括号、右括号或空，使未匹配左括号数量变成一个可能区间。暴力版的慢点要落到本题：慢版本会围绕“星号可以作为左括号、右括号或空，使未匹配左括号数量变成一个可能区间”枚举选择顺序。样例规律是“规律是左括号让区间加一，右括号让区间减一，星号让下界减一上界加一，且下界不能低于零”，说明存在一个局部选择或边界状态可以安全保留；不能证明这个选择不变差时，就不能把它叫贪心。后面的优化只消掉这类重复工作，不能改变答案集合。

\textbf{优化条件。} 本题的关键观察是：贪心维护可能未匹配左括号数的下界和上界，遇到星号扩大区间并把下界截到零。这句话必须能翻译成可维护状态；如果题目条件改变后观察不再成立，就不能继续套原来的写法。

\textbf{相邻状态。} 规律是左括号让区间加一，右括号让区间减一，星号让下界减一上界加一，且下界不能低于零。把这个特例继续拆开看：先构造一个非贪心选择，再比较两者留下的资源、右边界或剩余时间。只有能把非贪心第一步交换成贪心第一步且不变差，局部选择才是规律。

\textbf{状态复用。} 把指数级选择压成一个可证明安全的局部选择；状态通常是当前最远边界、最早结束时间、最小代价或剩余资源。本题落点是：规律是左括号让区间加一，右括号让区间减一，星号让下界减一上界加一，且下界不能低于零。手算时只改被新输入影响的那一小部分，而不是回头重算完整候选。

\textbf{指针和字段。} 状态字段要能说清：处理顺序、当前已选方案、剩余资源或最远边界、答案计数；若使用排序，要说明排序键；每次局部选择后要能说明当前状态不落后。手算时按这个协议跟踪：拿 (*) 手算，星号可以当空或左括号；拿 (*))，星号必须当左括号才能匹配后面的右括号。与其枚举星号三种选择，不如维护可能未匹配左括号数的区间 [low, high]。然后按协议复查：手算时同时写一个非贪心选择，与贪心选择比较剩余资源。若无法说明贪心后剩余资源不差，就不能直接下结论。

\textbf{代码落点。} 把局部规则写成业务含义；若需要排序，就处理相同关键字；循环中只维护证明需要的状态，不把局部选择和答案更新混在一起。关键代码行必须能逐句翻译成“旧状态怎样变成新状态”，否则就只是模板。

\textbf{正确性。} 证明从交换或领先性开始：任意最优解都能替换成当前贪心选择而不变差，或贪心状态始终不落后。

\textbf{边界实验。} 先用上界变负、最终下界不为零、连续星号、前缀右括号过多做反例检查。实验时除了上界变负、最终下界不为零、连续星号、前缀右括号过多，还要故意替换成一个看似合理但错误的局部规则，构造反例。能解释错误贪心为何失败，才算理解正确贪心。

\textbf{复杂度变化。} 暴力枚举选择顺序可能指数级；贪心通常先排序，再线性扫描或用堆维护少量候选。

\textbf{迁移追问。} 如果题目改成“也可用双栈保存左括号和星号位置，但区间贪心更能说明状态压缩”，先判断原状态是否仍然覆盖新答案集合，再决定是微调边界、改 value 语义，还是换算法。

\noindent\textbf{LeetCode 721 账户合并。}

\textbf{题目说明。} LeetCode 721 账户合并 的题面入口要先还原成具体任务：共享邮箱说明账户属于同一人，邮箱之间形成连通分量。

\textbf{样例入口。} 拿 John 的账号 [a, b] 和另一个 [b, c] 手算，邮箱 b 重叠，所以 a、b、c 属于同一人；名字只是输出标签，合并依据是邮箱。

\textbf{暴力基线。} 慢版本每次合并后用 DFS 或 BFS 重新判断连通性。它正确但重复遍历已有连接；当关系只会合并不会拆开时，可以压成集合代表。放到本题就是：先按“共享邮箱说明账户属于同一人，邮箱之间形成连通分量”完整试慢版本；样例里已经能看到“规律是邮箱映射到编号后并查集合并，最后按代表收集并排序邮箱”，所以优化前必须先能说清慢版本枚举了哪些候选。

\textbf{暴力过程。} 拿 John 的账号 [a, b] 和另一个 [b, c] 手算，邮箱 b 重叠，所以 a、b、c 属于同一人；名字只是输出标签，合并依据是邮箱。

\textbf{重复工作。} 题面模型：共享邮箱说明账户属于同一人，邮箱之间形成连通分量。暴力版的慢点要落到本题：慢版本会围绕“共享邮箱说明账户属于同一人，邮箱之间形成连通分量”反复搜索连通性。样例规律是“规律是邮箱映射到编号后并查集合并，最后按代表收集并排序邮箱”，说明关系只需要被压成集合代表；每次 union 失败或成功都要能翻译回题意。后面的优化只消掉这类重复工作，不能改变答案集合。

\textbf{优化条件。} 本题的关键观察是：给邮箱编号并查集合并，同根邮箱收集后排序输出。这句话必须能翻译成可维护状态；如果题目条件改变后观察不再成立，就不能继续套原来的写法。

\textbf{相邻状态。} 规律是邮箱映射到编号后并查集合并，最后按代表收集并排序邮箱。把这个特例继续拆开看：union 只是在维护等价类，不是在保存具体路径。两个节点代表元相同只能说明连通或关系一致；如果题目还要距离、比例或奇偶性，必须把附加信息挂到根关系上。

\textbf{状态复用。} 把连通分量压成代表元；查询只比较代表，合并只发生在两个根之间。本题落点是：规律是邮箱映射到编号后并查集合并，最后按代表收集并排序邮箱。手算时只改被新输入影响的那一小部分，而不是回头重算完整候选。

\textbf{指针和字段。} 状态字段要能说清：parent、size 或 rank、find 返回的代表元、合并时的附加信息；带权或奇偶并查集还要维护到根的关系。手算时按这个协议跟踪：拿 John 的账号 [a, b] 和另一个 [b, c] 手算，邮箱 b 重叠，所以 a、b、c 属于同一人；名字只是输出标签，合并依据是邮箱。然后按协议复查：手算时写 parent 表和集合代表。每次 union 只改变根之间的关系，find 后代表相同才说明连通。

\textbf{代码落点。} find 路径压缩不能破坏附加权值；union 只能连接两个根；合并失败和重复边要保留题意解释。关键代码行必须能逐句翻译成“旧状态怎样变成新状态”，否则就只是模板。

\textbf{正确性。} 证明从等价关系开始：合并维护连通性的传递性，find 返回的代表元就是当前集合身份。

\textbf{边界实验。} 先用同名不同人、一个账户多个邮箱、重复邮箱、输出排序做反例检查。实验时覆盖同名不同人、一个账户多个邮箱、重复邮箱、输出排序，打印每个元素的代表元和集合大小。重复合并、已经连通的边和字符串编号最容易暴露实现问题。

\textbf{复杂度变化。} 暴力每次查询都搜索连通性代价高；并查集把合并和查询摊还到近似常数，整体接近线性。

\textbf{迁移追问。} 如果题目改成“姓名不能作为合并依据，只能作为输出字段”，先判断原状态是否仍然覆盖新答案集合，再决定是微调边界、改 value 语义，还是换算法。

\noindent\textbf{LeetCode 815 公交路线。}

\textbf{题目说明。} LeetCode 815 公交路线 的题面入口要先还原成具体任务：公交路线和站点构成二分关系，换乘次数比站点步数更重要。

\textbf{样例入口。} 拿起点站可坐路线 A，A 和 B 在某站相交，B 到终点手算。按站点一步步 BFS 会反复展开同一条路线的所有站；按路线 BFS，一次乘车层数就是换乘次数。

\textbf{暴力基线。} 慢版本先按题意画出完整状态图，用普通搜索跑通可达性。这个过程用于确认公交路线和站点构成二分关系，换乘次数比站点步数更重要的节点和边，之后再讨论访问标记、层次或代价顺序。放到本题就是：先按“公交路线和站点构成二分关系，换乘次数比站点步数更重要”完整试慢版本；样例里已经能看到“规律是访问过的路线不再展开，站点到路线表只用于找到下一批路线”，所以优化前必须先能说清慢版本枚举了哪些候选。

\textbf{暴力过程。} 拿起点站可坐路线 A，A 和 B 在某站相交，B 到终点手算。按站点一步步 BFS 会反复展开同一条路线的所有站；按路线 BFS，一次乘车层数就是换乘次数。

\textbf{重复工作。} 题面模型：公交路线和站点构成二分关系，换乘次数比站点步数更重要。暴力版的慢点要落到本题：慢版本会围绕“公交路线和站点构成二分关系，换乘次数比站点步数更重要”反复展开同一批状态。样例规律是“规律是访问过的路线不再展开，站点到路线表只用于找到下一批路线”，说明必须把节点、边、访问标记和资源维度一次建清；同一状态不应重复入队，不同资源状态也不能误合并。后面的优化只消掉这类重复工作，不能改变答案集合。

\textbf{优化条件。} 本题的关键观察是：用站点到路线列表的哈希表找可乘路线，以路线或乘车层数做 BFS，访问过的路线不再重复展开。这句话必须能翻译成可维护状态；如果题目条件改变后观察不再成立，就不能继续套原来的写法。

\textbf{相邻状态。} 规律是访问过的路线不再展开，站点到路线表只用于找到下一批路线。把这个特例继续拆开看：状态节点不一定等于原始位置，还可能包含钥匙、剩余资源、时间或访问集合。visited 只能合并未来能力完全相同的状态，否则会把合法路径剪掉。

\textbf{状态复用。} 把重复搜索压成访问标记、距离数组或拓扑状态；邻居生成也要从全量扫描变成结构化枚举。本题落点是：规律是访问过的路线不再展开，站点到路线表只用于找到下一批路线。手算时只改被新输入影响的那一小部分，而不是回头重算完整候选。

\textbf{指针和字段。} 状态字段要能说清：状态编号、邻接生成方式、访问标记、距离或层数、队列内容；若状态含资源，visited 不能只按位置记录。手算时按这个协议跟踪：拿起点站可坐路线 A，A 和 B 在某站相交，B 到终点手算。按站点一步步 BFS 会反复展开同一条路线的所有站；按路线 BFS，一次乘车层数就是换乘次数。然后按协议复查：手算时画队列或栈，状态必须包含题目需要的所有资源。入队时标记还是出队时标记，要和去重语义一致。

\textbf{代码落点。} 入队时标记还是出队时标记必须一致；方向数组集中定义；多源 BFS 要一次性把所有源点放入初始队列。关键代码行必须能逐句翻译成“旧状态怎样变成新状态”，否则就只是模板。

\textbf{正确性。} 证明从可达性开始：搜索覆盖所有合法边能到达的状态，访问标记只删除重复状态而不删除不同未来能力。

\textbf{边界实验。} 先用起点等于终点、某站经过路线很多、路线重复访问、目标站不可达做反例检查。实验时覆盖起点等于终点、某站经过路线很多、路线重复访问、目标站不可达，并打印入队时的完整状态。若同一位置但资源不同的状态被合并，很多图题会出现隐蔽错误。

\textbf{复杂度变化。} 暴力重复从多个起点搜索会反复遍历图；正确建模和访问标记后，BFS 或 DFS 通常按节点数加边数计算，状态图题则按完整状态数计算。

\textbf{迁移追问。} 如果题目改成“若把每个站点之间完全建边会爆炸，应利用路线这一层压缩邻居生成”，先判断原状态是否仍然覆盖新答案集合，再决定是微调边界、改 value 语义，还是换算法。

\noindent\textbf{LeetCode 862 和至少为 K 的最短子数组。}

\textbf{题目说明。} LeetCode 862 和至少为 K 的最短子数组 的题面入口要先还原成具体任务：数组允许负数，普通滑动窗口的和不再随左端移动单调变化。

\textbf{样例入口。} 拿 [2, -1, 2]、K=3 手算，普通正数窗口会被 -1 破坏单调性；前缀和为 0, 2, 1, 3，只有 3-0 达到 3。

\textbf{暴力基线。} 暴力枚举所有 left/right，用前缀和 O(1) 求区间和，检查是否至少为 K 并更新最短长度。

\textbf{暴力过程。} 以 [2, -1, 2]、K=3 为例，前缀和是 [0, 2, 1, 3]。普通正数窗口会被 -1 破坏：右端增加后和可能变小。

\textbf{重复工作。} 题面模型：数组允许负数，普通滑动窗口的和不再随左端移动单调变化。暴力版的慢点要落到本题：浪费在每个右端都回头扫描所有历史左端。需要的是某些前缀和尽量小且下标尽量靠后的候选左端。后面的优化只消掉这类重复工作，不能改变答案集合。

\textbf{优化条件。} 本题的关键观察是：在前缀和数组上用单调队列维护可能成为最优左端的下标。这句话必须能翻译成可维护状态；如果题目条件改变后观察不再成立，就不能继续套原来的写法。

\textbf{相邻状态。} 规律是在前缀和上用单调队列维护可能的最优左端，队首能满足时结算长度。把这个特例继续拆开看：右端进入只带来一个新元素，左端离开只删掉一个旧元素；只有当旧左端被移走后不可能再产生更优答案时，窗口才可以单向推进。若这个单调淘汰条件不存在，就要换成前缀和、队列或其他状态。

\textbf{状态复用。} 构造前缀和 prefix。用单调递增队列保存可能作为左端的前缀下标；当前前缀减队首前缀 >=K 时更新答案并弹出队首；队尾前缀不小于当前前缀时被当前支配，弹出。手算时只改被新输入影响的那一小部分，而不是回头重算完整候选。

\textbf{指针和字段。} 状态字段要能说清：deque 中保存前缀下标，且对应 prefix 值递增；i 是当前右端后一位，i-deque.front 是候选长度。手算时按这个协议跟踪：prefix 到 3 时，3-0>=3，长度 3；若出现更小的新前缀，它会支配队尾较大前缀，因为未来用它能得到更大区间和且更短或不更差。

\textbf{代码落点。} 关键有两个 while：先从队首结算满足 K 的最短长度，再从队尾删除 prefix 值不小于当前值的被支配候选。关键代码行必须能逐句翻译成“旧状态怎样变成新状态”，否则就只是模板。

\textbf{正确性。} 队首是当前最早且前缀较小的候选；一旦满足 K，继续保留它只会让未来长度更长，所以可以弹出。队尾被更小且更新的前缀支配，永远不会更优。

\textbf{边界实验。} 先用负数、答案不存在、K 为负或零、队尾支配关系做反例检查。实验时重点跑负数、答案不存在、K 为负或零、队尾支配关系，并打印左右端、窗口计数和答案。若左端发生回退，或者窗口失去合法性后仍更新答案，就能很快定位错误。

\textbf{复杂度变化。} 暴力 O(\ensuremath{n^{2}})，单调队列中每个前缀入队出队各一次，时间 O(n)，空间 O(n)。

\textbf{迁移追问。} 如果题目改成“这是从窗口题升级到前缀和加单调队列的代表”，先判断原状态是否仍然覆盖新答案集合，再决定是微调边界、改 value 语义，还是换算法。

\noindent\textbf{LeetCode 895 最大频率栈。}

\textbf{题目说明。} LeetCode 895 最大频率栈 的题面入口要先还原成具体任务：弹出时既要频率最高，又要在同频中最近入栈。

\textbf{样例入口。} 拿 push 5, 7, 5, 7, 4, 5 后 pop 手算，频率最高是 5，先弹 5；下一次 5 和 7 同频，弹最近达到该频率的 7。

\textbf{暴力基线。} 慢版本让每个位置向左或向右线性寻找边界，或者让每个结构反复等待匹配。它能算出答案，但很多尚未完成的候选会被未来同一个事件一起解决。放到本题就是：先按“弹出时既要频率最高，又要在同频中最近入栈”完整试慢版本；样例里已经能看到“规律是值到频率、频率到栈、当前最大频率三者共同维护”，所以优化前必须先能说清慢版本枚举了哪些候选。

\textbf{暴力过程。} 拿 push 5, 7, 5, 7, 4, 5 后 pop 手算，频率最高是 5，先弹 5；下一次 5 和 7 同频，弹最近达到该频率的 7。

\textbf{重复工作。} 题面模型：弹出时既要频率最高，又要在同频中最近入栈。暴力版的慢点要落到本题：慢版本会围绕“弹出时既要频率最高，又要在同频中最近入栈”反复寻找最近边界或最近未完成对象。样例规律是“规律是值到频率、频率到栈、当前最大频率三者共同维护”，说明旧候选一旦遇到当前输入就能被批量结算；继续为每个位置单独向左或向右扫描就是浪费。后面的优化只消掉这类重复工作，不能改变答案集合。

\textbf{优化条件。} 本题的关键观察是：哈希表记录值频率，频率到栈记录达到该频率的入栈顺序，维护当前最大频率。这句话必须能翻译成可维护状态；如果题目条件改变后观察不再成立，就不能继续套原来的写法。

\textbf{相邻状态。} 规律是值到频率、频率到栈、当前最大频率三者共同维护。把这个特例继续拆开看：栈里的对象都是答案尚未确定的候选；一旦当前元素触发弹栈，被弹出的候选必须已经同时拿到了左边界和右边界，未来元素不再有资格改变它的答案。

\textbf{状态复用。} 把反复寻找最近边界改成维护未完成候选；新元素只和栈顶比较，连续弹出一批已经被当前元素解决的对象。本题落点是：规律是值到频率、频率到栈、当前最大频率三者共同维护。手算时只改被新输入影响的那一小部分，而不是回头重算完整候选。

\textbf{指针和字段。} 状态字段要能说清：栈内对象的业务含义、当前输入、弹出时结算的对象、左边界和右边界；栈中存下标还是值要明确。手算时按这个协议跟踪：拿 push 5, 7, 5, 7, 4, 5 后 pop 手算，频率最高是 5，先弹 5；下一次 5 和 7 同频，弹最近达到该频率的 7。然后按协议复查：手算时记录栈内元素的业务含义，不只记录数值。入栈表示答案未定，弹栈表示当前事件已经让它的答案定下来。

\textbf{代码落点。} 弹栈前确认栈非空，弹出后再读取新的左边界；相等元素和哨兵高度要有统一策略。关键代码行必须能逐句翻译成“旧状态怎样变成新状态”，否则就只是模板。

\textbf{正确性。} 证明从弹出时机开始：被弹出的对象在当前时刻答案已经确定，未来元素无法再改变它。

\textbf{边界实验。} 先用重复 push、pop 后最大频率下降、同频最近性、空栈不在题意内做反例检查。实验时准备重复 push、pop 后最大频率下降、同频最近性、空栈不在题意内，打印每次入栈、弹栈和结算对象。重点观察相等元素、空栈、哨兵和最后一次结算。

\textbf{复杂度变化。} 暴力为每个位置寻找边界常见为平方级；单调栈让每个元素最多入栈一次、出栈一次，因此主体扫描为线性时间。

\textbf{迁移追问。} 如果题目改成“它是栈和哈希表按频率维度组合的设计题”，先判断原状态是否仍然覆盖新答案集合，再决定是微调边界、改 value 语义，还是换算法。

\noindent\textbf{LeetCode 981 基于时间的键值存储。}

\textbf{题目说明。} LeetCode 981 基于时间的键值存储 的题面入口要先还原成具体任务：同一个 key 下的记录按时间戳递增，查询是不超过目标时间的最新值。

\textbf{样例入口。} 拿 key=foo 的记录 (1, bar)、(4, baz)，查询 t=3 手算，答案应是时间不超过 3 的最后一条，也就是 bar。查询 t=4 才是 baz。

\textbf{暴力基线。} 慢版本按顺序扫描下标、逐个试答案值，或直接检查候选直到遇到目标边界。它先保证候选空间覆盖完整，但每次只排除一个候选，浪费了题目给出的顺序、比较或可行性边界。放到本题就是：先按“同一个 key 下的记录按时间戳递增，查询是不超过目标时间的最新值”完整试慢版本；样例里已经能看到“规律是每个 key 内部时间递增，二分第一个大于查询时间的位置，再回退一格”，所以优化前必须先能说清慢版本枚举了哪些候选。

\textbf{暴力过程。} 拿 key=foo 的记录 (1, bar)、(4, baz)，查询 t=3 手算，答案应是时间不超过 3 的最后一条，也就是 bar。查询 t=4 才是 baz。

\textbf{重复工作。} 题面模型：同一个 key 下的记录按时间戳递增，查询是不超过目标时间的最新值。暴力版的慢点要落到本题：慢版本会按顺序试“同一个 key 下的记录按时间戳递增，查询是不超过目标时间的最新值”里的候选；而样例规律已经指出“规律是每个 key 内部时间递增，二分第一个大于查询时间的位置，再回退一格”。一次比较或一次可行性检查能丢掉一整段候选，继续线性试就是浪费。后面的优化只消掉这类重复工作，不能改变答案集合。

\textbf{优化条件。} 本题的关键观察是：哈希表定位 key 后，在该 key 的时间列表中二分第一个大于目标时间的位置并回退一格。这句话必须能翻译成可维护状态；如果题目条件改变后观察不再成立，就不能继续套原来的写法。

\textbf{相邻状态。} 规律是每个 key 内部时间递增，二分第一个大于查询时间的位置，再回退一格。把这个特例继续拆开看：每次比较 mid 之后，必须能指出哪一侧整体不可能包含目标，或哪一侧整体已经满足可行性。二分的核心不是 mid 公式，而是被丢弃区间确实不含答案。

\textbf{状态复用。} 把逐个试候选改成维护答案所在区间；边界谓词、可行性检查或局部比较给出分支方向，每个分支都必须能排除一侧候选。本题落点是：规律是每个 key 内部时间递增，二分第一个大于查询时间的位置，再回退一格。手算时只改被新输入影响的那一小部分，而不是回头重算完整候选。

\textbf{指针和字段。} 状态字段要能说清：left、right、mid、mid 处比较值或检查返回值、答案区间含义、返回值语义；每一轮更新都要保留目标位置或第一个可行值。手算时按这个协议跟踪：拿 key=foo 的记录 (1, bar)、(4, baz)，查询 t=3 手算，答案应是时间不超过 3 的最后一条，也就是 bar。查询 t=4 才是 baz。然后按协议复查：手算时列出 left、mid、right，以及 mid 处比较结果、可行性检查结果或局部趋势。每一步都要写出被丢弃的一侧为什么不含答案。

\textbf{代码落点。} 检查函数或比较分支单独写清，mid 不溢出，循环退出条件和返回值配套；每个分支都要解释丢弃区间。关键代码行必须能逐句翻译成“旧状态怎样变成新状态”，否则就只是模板。

\textbf{正确性。} 证明从排除一侧开始：答案空间题证明谓词单调，局部比较题证明保留下来的区间仍含目标；不能只靠经验猜测左右边界。

\textbf{边界实验。} 先用key 不存在、查询早于第一条记录、相同 key 多次 set、时间戳递增假设做反例检查。实验时覆盖key 不存在、查询早于第一条记录、相同 key 多次 set、时间戳递增假设，逐轮打印 left、mid、right、比较值或检查结果。只要有一次被排除区间仍含答案，二分不变量就错了。

\textbf{复杂度变化。} 暴力逐个试候选是线性或和值域成正比；二分每次排除一半候选，复杂度变成对数级，答案二分还要乘上单次检查成本。

\textbf{迁移追问。} 如果题目改成“如果时间戳无序到达，必须先排序或改成有序表维护每个 key 的记录”，先判断原状态是否仍然覆盖新答案集合，再决定是微调边界、改 value 语义，还是换算法。

\noindent\textbf{LeetCode 1091 二进制矩阵中的最短路径。}

\textbf{题目说明。} LeetCode 1091 二进制矩阵中的最短路径 的题面入口要先还原成具体任务：八方向网格无权边，答案是从左上到右下的最少格子数。

\textbf{样例入口。} 拿 [[0, 1], [1, 0]] 手算，从左上可以斜着一步到右下，路径长度是 2；若起点或终点为 1 直接不可达。

\textbf{暴力基线。} 慢版本先按题意画出完整状态图，用普通搜索跑通可达性。这个过程用于确认八方向网格无权边，答案是从左上到右下的最少格子数的节点和边，之后再讨论访问标记、层次或代价顺序。放到本题就是：先按“八方向网格无权边，答案是从左上到右下的最少格子数”完整试慢版本；样例里已经能看到“规律是八方向 BFS，第一次到达终点就是最短路径长度”，所以优化前必须先能说清慢版本枚举了哪些候选。

\textbf{暴力过程。} 拿 [[0, 1], [1, 0]] 手算，从左上可以斜着一步到右下，路径长度是 2；若起点或终点为 1 直接不可达。

\textbf{重复工作。} 题面模型：八方向网格无权边，答案是从左上到右下的最少格子数。暴力版的慢点要落到本题：慢版本会围绕“八方向网格无权边，答案是从左上到右下的最少格子数”反复展开同一批状态。样例规律是“规律是八方向 BFS，第一次到达终点就是最短路径长度”，说明必须把节点、边、访问标记和资源维度一次建清；同一状态不应重复入队，不同资源状态也不能误合并。后面的优化只消掉这类重复工作，不能改变答案集合。

\textbf{优化条件。} 本题的关键观察是：BFS 入队时标记访问，层数或距离随状态保存。这句话必须能翻译成可维护状态；如果题目条件改变后观察不再成立，就不能继续套原来的写法。

\textbf{相邻状态。} 规律是八方向 BFS，第一次到达终点就是最短路径长度。把这个特例继续拆开看：状态节点不一定等于原始位置，还可能包含钥匙、剩余资源、时间或访问集合。visited 只能合并未来能力完全相同的状态，否则会把合法路径剪掉。

\textbf{状态复用。} 把重复搜索压成访问标记、距离数组或拓扑状态；邻居生成也要从全量扫描变成结构化枚举。本题落点是：规律是八方向 BFS，第一次到达终点就是最短路径长度。手算时只改被新输入影响的那一小部分，而不是回头重算完整候选。

\textbf{指针和字段。} 状态字段要能说清：状态编号、邻接生成方式、访问标记、距离或层数、队列内容；若状态含资源，visited 不能只按位置记录。手算时按这个协议跟踪：拿 [[0, 1], [1, 0]] 手算，从左上可以斜着一步到右下，路径长度是 2；若起点或终点为 1 直接不可达。然后按协议复查：手算时画队列或栈，状态必须包含题目需要的所有资源。入队时标记还是出队时标记，要和去重语义一致。

\textbf{代码落点。} 入队时标记还是出队时标记必须一致；方向数组集中定义；多源 BFS 要一次性把所有源点放入初始队列。关键代码行必须能逐句翻译成“旧状态怎样变成新状态”，否则就只是模板。

\textbf{正确性。} 证明从可达性开始：搜索覆盖所有合法边能到达的状态，访问标记只删除重复状态而不删除不同未来能力。

\textbf{边界实验。} 先用起点或终点阻塞、单格、八方向越界、无路可达做反例检查。实验时覆盖起点或终点阻塞、单格、八方向越界、无路可达，并打印入队时的完整状态。若同一位置但资源不同的状态被合并，很多图题会出现隐蔽错误。

\textbf{复杂度变化。} 暴力重复从多个起点搜索会反复遍历图；正确建模和访问标记后，BFS 或 DFS 通常按节点数加边数计算，状态图题则按完整状态数计算。

\textbf{迁移追问。} 如果题目改成“边权相同才能用 BFS，若每格有代价就要最短路”，先判断原状态是否仍然覆盖新答案集合，再决定是微调边界、改 value 语义，还是换算法。

\noindent\textbf{LeetCode 1235 规划兼职工作。}

\textbf{题目说明。} LeetCode 1235 规划兼职工作 的题面入口要先还原成具体任务：每份工作有开始、结束和收益，选择不冲突工作使收益最大。

\textbf{样例入口。} 拿 jobs (1, 3, 50), (2, 4, 10), (3, 5, 40), (3, 6, 70) 手算，选 (1, 3, 50) 后还能接 (3, 6, 70)，总 120。

\textbf{暴力基线。} 慢版本先写递归搜索树：节点表示处理位置、剩余资源和当前选择。只要不同路径反复到达同一个节点，就说明需要把这个节点命名成 DP 状态。放到本题就是：先按“每份工作有开始、结束和收益，选择不冲突工作使收益最大”完整试慢版本；样例里已经能看到“规律是按结束时间排序，dp[i] 在不选当前和选当前加上前一个不冲突工作之间取最大”，所以优化前必须先能说清慢版本枚举了哪些候选。

\textbf{暴力过程。} 拿 jobs (1, 3, 50), (2, 4, 10), (3, 5, 40), (3, 6, 70) 手算，选 (1, 3, 50) 后还能接 (3, 6, 70)，总 120。

\textbf{重复工作。} 题面模型：每份工作有开始、结束和收益，选择不冲突工作使收益最大。暴力版的慢点要落到本题：慢版本会在“每份工作有开始、结束和收益，选择不冲突工作使收益最大”的选择树里反复到达相同参数。样例规律是“规律是按结束时间排序，dp[i] 在不选当前和选当前加上前一个不冲突工作之间取最大”，说明这个参数组合可以命名成状态，算过之后后面直接复用。后面的优化只消掉这类重复工作，不能改变答案集合。

\textbf{优化条件。} 本题的关键观察是：按结束时间排序，dp[i] 在不选当前和选当前加上兼容前驱之间取最大，前驱用二分找。这句话必须能翻译成可维护状态；如果题目条件改变后观察不再成立，就不能继续套原来的写法。

\textbf{相邻状态。} 规律是按结束时间排序，dp[i] 在不选当前和选当前加上前一个不冲突工作之间取最大。把这个特例继续拆开看：先问暴力搜索里哪些不同路径到达同一个后续问题，再给这个后续问题命名为状态。状态必须包含未来决策需要的全部信息，转移只从更小且已经算清的状态来。

\textbf{状态复用。} 把重复递归节点命名为状态；遍历顺序只是在保证依赖状态已经算完。本题落点是：规律是按结束时间排序，dp[i] 在不选当前和选当前加上前一个不冲突工作之间取最大。手算时只改被新输入影响的那一小部分，而不是回头重算完整候选。

\textbf{指针和字段。} 状态字段要能说清：状态下标、状态值语义、初始化、转移来源、遍历顺序；空间压缩时还要指出读到的是旧值还是新值。手算时按这个协议跟踪：拿 jobs (1, 3, 50), (2, 4, 10), (3, 5, 40), (3, 6, 70) 手算，选 (1, 3, 50) 后还能接 (3, 6, 70)，总 120。然后按协议复查：手算时先写一个很小的 DP 表或记忆化搜索记录。每个格子旁边写它来自哪些更小状态。

\textbf{代码落点。} 先写未压缩版本校验转移；压缩前后用小表对拍；不可达哨兵参与加法前要检查溢出。关键代码行必须能逐句翻译成“旧状态怎样变成新状态”，否则就只是模板。

\textbf{正确性。} 证明从最后一步开始：任意最优解的最后一步都在转移枚举中，转移来源已经由更小状态正确计算。

\textbf{边界实验。} 先用结束等于开始是否兼容、收益相同、工作排序、空列表做反例检查。实验时覆盖结束等于开始是否兼容、收益相同、工作排序、空列表，先打印完整 DP 表，再比较空间压缩版本。若压缩版不同，优先检查遍历顺序和旧值保存。

\textbf{复杂度变化。} 暴力递归会重复计算相同子问题，常见指数级；DP 把每个状态算一次，时间是状态数乘单个转移成本，空间是状态表大小或压缩后的大小。

\textbf{迁移追问。} 如果题目改成“这是排序、二分和 DP 的组合，不是按收益贪心”，先判断原状态是否仍然覆盖新答案集合，再决定是微调边界、改 value 语义，还是换算法。

\noindent\textbf{LeetCode 1248 统计优美子数组。}

\textbf{题目说明。} LeetCode 1248 统计优美子数组 的题面入口要先还原成具体任务：恰好 K 个奇数可转成前缀奇数计数差为 K。

\textbf{样例入口。} 拿 [1, 1, 2, 1, 1]、k=3 手算，奇数个数前缀从 0 到 3 时，需要历史奇数个数 0；每个匹配历史前缀对应一个优美子数组。

\textbf{暴力基线。} 慢版本按题意两两比较、枚举区间、扫描历史或持续迭代并查重。把检查条件写出来后，通常能看到当前状态只缺一个历史状态、规范键或已经访问过的中间状态；这正是从平方枚举或重复迭代走向哈希表的入口。放到本题就是：先按“恰好 K 个奇数可转成前缀奇数计数差为 K”完整试慢版本；样例里已经能看到“规律是把奇数看成 1、偶数看成 0，转成前缀和计数”，所以优化前必须先能说清慢版本枚举了哪些候选。

\textbf{暴力过程。} 拿 [1, 1, 2, 1, 1]、k=3 手算，奇数个数前缀从 0 到 3 时，需要历史奇数个数 0；每个匹配历史前缀对应一个优美子数组。

\textbf{重复工作。} 题面模型：恰好 K 个奇数可转成前缀奇数计数差为 K。暴力版的慢点要落到本题：慢版本会为“恰好 K 个奇数可转成前缀奇数计数差为 K”反复寻找历史状态；而样例规律已经指出“规律是把奇数看成 1、偶数看成 0，转成前缀和计数”。这些补数、前缀状态、规范键、半边和或访问过的中间状态可以用一个 key 归类，不必每次重新扫描或重复比较。后面的优化只消掉这类重复工作，不能改变答案集合。

\textbf{优化条件。} 本题的关键观察是：哈希表保存某个奇数计数出现次数，当前计数查询 count - K。这句话必须能翻译成可维护状态；如果题目条件改变后观察不再成立，就不能继续套原来的写法。

\textbf{相邻状态。} 规律是把奇数看成 1、偶数看成 0，转成前缀和计数。把这个特例继续拆开看：每个合法答案都要还原成当前状态和某个历史状态的配对；哈希表的 key 负责定位历史状态，value 负责表达次数、最早位置或存在性。先查还是先写，必须由是否允许当前前缀和自己配对来决定。

\textbf{状态复用。} 把回头扫描、两两比较或重复状态检测改成查表、分桶或访问记录；key 是当前题目需要的补数、前缀状态、规范表示、半边和或中间状态，value 根据目标选择次数、下标、列表或存在性。本题落点是：规律是把奇数看成 1、偶数看成 0，转成前缀和计数。手算时只改被新输入影响的那一小部分，而不是回头重算完整候选。

\textbf{指针和字段。} 状态字段要能说清：当前状态或规范键、需要查询的历史 key、哈希表 value 语义、答案累计值或分桶结果；空前缀、空字符串、重复状态或初始状态要单独说明。手算时按这个协议跟踪：拿 [1, 1, 2, 1, 1]、k=3 手算，奇数个数前缀从 0 到 3 时，需要历史奇数个数 0；每个匹配历史前缀对应一个优美子数组。然后按协议复查：手算时列出当前输入、当前状态或规范键、需要查询的历史 key、查到的 value、更新后的哈希表。这样能看出先查后插、先分桶后输出，还是先检测重复状态。

\textbf{代码落点。} 先查询还是先插入要由题意决定；key 构造必须无歧义；哈希表 value 的默认插入副作用要可控。关键代码行必须能逐句翻译成“旧状态怎样变成新状态”，否则就只是模板。

\textbf{正确性。} 证明从状态等价开始：任意合法答案都对应当前状态和某个历史 key、规范键或访问状态，查到的 key 也能反推出合法答案或合法分组。

\textbf{边界实验。} 先用K 为零、全偶数、全奇数、从开头开始的区间做反例检查。实验时把K 为零、全偶数、全奇数、从开头开始的区间加入对拍集合，用平方暴力和优化版比较。失败时先看初始历史状态、查询更新顺序和哈希表 value 类型。

\textbf{复杂度变化。} 暴力两两比较、枚举区间、扫描历史或重复迭代常见为平方级或更多；把历史状态、规范键或访问状态放进哈希表后，每步只做常数次查询和更新，通常是线性时间、线性空间。

\textbf{迁移追问。} 如果题目改成“也可用至多 K 减至多 K-1 的滑动窗口”，先判断原状态是否仍然覆盖新答案集合，再决定是微调边界、改 value 语义，还是换算法。

\noindent\textbf{LeetCode 1438 绝对差不超过限制的最长连续子数组。}

\textbf{题目说明。} LeetCode 1438 绝对差不超过限制的最长连续子数组 的题面入口要先还原成具体任务：窗口合法性取决于窗口内最大值和最小值之差。

\textbf{样例入口。} 拿 [8, 2, 4, 7]、limit=4 手算，窗口 [8, 2] 最大最小差 6，不合法，左端必须右移；窗口 [2, 4] 合法。

\textbf{暴力基线。} 暴力枚举每个窗口，重新扫描窗口求最大值和最小值，若 max-min<=limit 就更新最长长度。

\textbf{暴力过程。} 以 [8, 2, 4, 7]、limit=4 为例，窗口 [8, 2] 差为 6 不合法，窗口 [2, 4] 合法。暴力每次都重新找最大最小。

\textbf{重复工作。} 题面模型：窗口合法性取决于窗口内最大值和最小值之差。暴力版的慢点要落到本题：浪费在相邻窗口重复求极值。右端新增一个元素后，只有被它支配的旧极值候选需要删除；左端移动时只需清理过期下标。后面的优化只消掉这类重复工作，不能改变答案集合。

\textbf{优化条件。} 本题的关键观察是：两个单调队列分别维护最大和最小，超过限制时移动左端并清理过期下标。这句话必须能翻译成可维护状态；如果题目条件改变后观察不再成立，就不能继续套原来的写法。

\textbf{相邻状态。} 规律是两个单调队列分别维护最大值和最小值，差值超限时移动左端并清理过期下标。把这个特例继续拆开看：右端进入只带来一个新元素，左端离开只删掉一个旧元素；只有当旧左端被移走后不可能再产生更优答案时，窗口才可以单向推进。若这个单调淘汰条件不存在，就要换成前缀和、队列或其他状态。

\textbf{状态复用。} 用 maxDeque 维护递减下标队列，队首是窗口最大值；minDeque 维护递增下标队列，队首是窗口最小值。差值超限时移动 left 并清理过期。手算时只改被新输入影响的那一小部分，而不是回头重算完整候选。

\textbf{指针和字段。} 状态字段要能说清：left/right 是窗口边界，maxDeque.front/minDeque.front 是当前极值下标，answer 是最长合法窗口长度。手算时按这个协议跟踪：加入 2 后，最大队首为 8、最小队首为 2，差 6 超限；left 从 0 移到 1，过期的 8 被清理，窗口 [2] 合法。

\textbf{代码落点。} 关键是入队前从队尾弹出被新值支配的下标；收缩时如果队首下标小于 left 就弹出。关键代码行必须能逐句翻译成“旧状态怎样变成新状态”，否则就只是模板。

\textbf{正确性。} 单调队列保留的都是可能成为当前或未来极值的下标，被新元素支配的旧下标既更旧又不更优，永远不会再成为极值。

\textbf{边界实验。} 先用limit 为零、重复元素、最大最小同时过期、窗口收缩多次做反例检查。实验时重点跑limit 为零、重复元素、最大最小同时过期、窗口收缩多次，并打印左右端、窗口计数和答案。若左端发生回退，或者窗口失去合法性后仍更新答案，就能很快定位错误。

\textbf{复杂度变化。} 暴力 O(\ensuremath{n^{2}})，两个单调队列每个下标入队出队各一次，时间 O(n)，空间 O(n)。

\textbf{迁移追问。} 如果题目改成“它说明窗口状态可以是动态极值，不只是频次或总和”，先判断原状态是否仍然覆盖新答案集合，再决定是微调边界、改 value 语义，还是换算法。

\noindent\textbf{LeetCode 1462 课程表 IV。}

\textbf{题目说明。} LeetCode 1462 课程表 IV 的题面入口要先还原成具体任务：大量先修关系查询需要预处理课程之间的可达性。

\textbf{样例入口。} 拿 0->1->2 和查询 0 是否先修 2 手算，直接边只能回答 0->1、1->2，不能回答传递关系。若查询很多，每次 DFS 会重复走相同路径。

\textbf{暴力基线。} 慢版本先按题意画出完整状态图，用普通搜索跑通可达性。这个过程用于确认大量先修关系查询需要预处理课程之间的可达性的节点和边，之后再讨论访问标记、层次或代价顺序。放到本题就是：先按“大量先修关系查询需要预处理课程之间的可达性”完整试慢版本；样例里已经能看到“规律是先做可达性闭包或拓扑合并祖先集合，再 O(1) 回答查询”，所以优化前必须先能说清慢版本枚举了哪些候选。

\textbf{暴力过程。} 拿 0->1->2 和查询 0 是否先修 2 手算，直接边只能回答 0->1、1->2，不能回答传递关系。若查询很多，每次 DFS 会重复走相同路径。

\textbf{重复工作。} 题面模型：大量先修关系查询需要预处理课程之间的可达性。暴力版的慢点要落到本题：慢版本会围绕“大量先修关系查询需要预处理课程之间的可达性”反复展开同一批状态。样例规律是“规律是先做可达性闭包或拓扑合并祖先集合，再 O(1) 回答查询”，说明必须把节点、边、访问标记和资源维度一次建清；同一状态不应重复入队，不同资源状态也不能误合并。后面的优化只消掉这类重复工作，不能改变答案集合。

\textbf{优化条件。} 本题的关键观察是：可以用 Floyd 传递闭包、每个点 BFS，或拓扑 DP 合并祖先集合，查询时直接查可达表。这句话必须能翻译成可维护状态；如果题目条件改变后观察不再成立，就不能继续套原来的写法。

\textbf{相邻状态。} 规律是先做可达性闭包或拓扑合并祖先集合，再 O(1) 回答查询。把这个特例继续拆开看：状态节点不一定等于原始位置，还可能包含钥匙、剩余资源、时间或访问集合。visited 只能合并未来能力完全相同的状态，否则会把合法路径剪掉。

\textbf{状态复用。} 把重复搜索压成访问标记、距离数组或拓扑状态；邻居生成也要从全量扫描变成结构化枚举。本题落点是：规律是先做可达性闭包或拓扑合并祖先集合，再 O(1) 回答查询。手算时只改被新输入影响的那一小部分，而不是回头重算完整候选。

\textbf{指针和字段。} 状态字段要能说清：状态编号、邻接生成方式、访问标记、距离或层数、队列内容；若状态含资源，visited 不能只按位置记录。手算时按这个协议跟踪：拿 0->1->2 和查询 0 是否先修 2 手算，直接边只能回答 0->1、1->2，不能回答传递关系。若查询很多，每次 DFS 会重复走相同路径。然后按协议复查：手算时画队列或栈，状态必须包含题目需要的所有资源。入队时标记还是出队时标记，要和去重语义一致。

\textbf{代码落点。} 入队时标记还是出队时标记必须一致；方向数组集中定义；多源 BFS 要一次性把所有源点放入初始队列。关键代码行必须能逐句翻译成“旧状态怎样变成新状态”，否则就只是模板。

\textbf{正确性。} 证明从可达性开始：搜索覆盖所有合法边能到达的状态，访问标记只删除重复状态而不删除不同未来能力。

\textbf{边界实验。} 先用课程之间没有关系、自身是否算先修、查询很多、图中按题意应为 DAG做反例检查。实验时覆盖课程之间没有关系、自身是否算先修、查询很多、图中按题意应为 DAG，并打印入队时的完整状态。若同一位置但资源不同的状态被合并，很多图题会出现隐蔽错误。

\textbf{复杂度变化。} 暴力重复从多个起点搜索会反复遍历图；正确建模和访问标记后，BFS 或 DFS 通常按节点数加边数计算，状态图题则按完整状态数计算。

\textbf{迁移追问。} 如果题目改成“如果查询很少，逐次搜索可能更省预处理成本；多查询时闭包更稳”，先判断原状态是否仍然覆盖新答案集合，再决定是微调边界、改 value 语义，还是换算法。

\noindent\textbf{LeetCode 1514 概率最大的路径。}

\textbf{题目说明。} LeetCode 1514 概率最大的路径 的题面入口要先还原成具体任务：路径概率是边概率乘积，目标是最大乘积路径。

\textbf{样例入口。} 拿 0->1 概率 0.5、1->2 概率 0.5、0->2 概率 0.2 手算，经 1 到 2 的概率 0.25 更大。

\textbf{暴力基线。} 慢版本可以枚举路径，或先用普通队列试跑一个小图。它会暴露步数最少和代价最小是否一致；一旦不一致，就必须围绕代价组织状态。放到本题就是：先按“路径概率是边概率乘积，目标是最大乘积路径”完整试慢版本；样例里已经能看到“规律是把距离改成最大概率，优先队列每次弹出当前概率最高的节点并做乘法松弛”，所以优化前必须先能说清慢版本枚举了哪些候选。

\textbf{暴力过程。} 拿 0->1 概率 0.5、1->2 概率 0.5、0->2 概率 0.2 手算，经 1 到 2 的概率 0.25 更大。

\textbf{重复工作。} 题面模型：路径概率是边概率乘积，目标是最大乘积路径。暴力版的慢点要落到本题：慢版本会围绕“路径概率是边概率乘积，目标是最大乘积路径”枚举路径或按错误顺序扩展状态。样例规律是“规律是把距离改成最大概率，优先队列每次弹出当前概率最高的节点并做乘法松弛”，说明队列或堆的弹出顺序必须和代价定义一致，旧状态为什么能跳过也要讲清。后面的优化只消掉这类重复工作，不能改变答案集合。

\textbf{优化条件。} 本题的关键观察是：用大顶堆每次扩展当前概率最高节点，乘以边概率得到候选概率。这句话必须能翻译成可维护状态；如果题目条件改变后观察不再成立，就不能继续套原来的写法。

\textbf{相邻状态。} 规律是把距离改成最大概率，优先队列每次弹出当前概率最高的节点并做乘法松弛。把这个特例继续拆开看：队列或堆弹出的顺序必须和代价规则一致；旧状态被跳过，是因为已经有更小代价到达同一状态。只要边权或代价合成方式改变，就要重新证明扩展顺序。

\textbf{状态复用。} 把路径枚举压成距离数组和优先队列；每次只扩展当前最有希望的未过期状态。本题落点是：规律是把距离改成最大概率，优先队列每次弹出当前概率最高的节点并做乘法松弛。手算时只改被新输入影响的那一小部分，而不是回头重算完整候选。

\textbf{指针和字段。} 状态字段要能说清：距离数组、优先队列状态、松弛候选、旧状态判定、不可达哨兵；资源约束题还要把资源维度放进状态。手算时按这个协议跟踪：拿 0->1 概率 0.5、1->2 概率 0.5、0->2 概率 0.2 手算，经 1 到 2 的概率 0.25 更大。然后按协议复查：手算时记录每次弹出的代价、当前距离数组和松弛结果。旧状态被弹出时为什么跳过，也要写在表里。

\textbf{代码落点。} 松弛时只在候选更优时更新；priority\_queue 弹出旧状态要跳过；距离加法用更宽整数。关键代码行必须能逐句翻译成“旧状态怎样变成新状态”，否则就只是模板。

\textbf{正确性。} 证明从扩展顺序开始：当前弹出的状态在代价规则下已经不会被未来路径改得更优。

\textbf{边界实验。} 先用不可达、概率为零、起点等于终点、浮点比较做反例检查。实验时覆盖不可达、概率为零、起点等于终点、浮点比较，打印每次弹出的代价和是否过期。若普通 BFS 与 Dijkstra 结果不同，要能解释边权条件造成的差异。

\textbf{复杂度变化。} 暴力枚举路径可能指数级；Dijkstra 在邻接表和堆实现下按边数、点数和堆操作计算，0-1 BFS 可按节点数加边数计算。

\textbf{迁移追问。} 如果题目改成“取负对数后可转成最短路，但直接最大堆更直观”，先判断原状态是否仍然覆盖新答案集合，再决定是微调边界、改 value 语义，还是换算法。

\noindent\textbf{LeetCode 1675 数组的最小偏移量。}

\textbf{题目说明。} LeetCode 1675 数组的最小偏移量 的题面入口要先还原成具体任务：奇数可以先乘二，偶数可以不断除二，目标是缩小当前最大最小差。

\textbf{样例入口。} 拿 [1, 2, 3, 4] 手算，奇数可先翻倍成偶数，之后只能不断把当前最大偶数减半；每次记录最大值和当前最小值差。

\textbf{暴力基线。} 慢版本每轮扫描全部候选来找最大或最小，再更新候选集合。它的答案选择过程清楚，但动态集合每变化一次就重新找极值，代价太高。放到本题就是：先按“奇数可以先乘二，偶数可以不断除二，目标是缩小当前最大最小差”完整试慢版本；样例里已经能看到“规律是把所有数变到可下降的最高状态，再用大顶堆逐步降低最大值”，所以优化前必须先能说清慢版本枚举了哪些候选。

\textbf{暴力过程。} 拿 [1, 2, 3, 4] 手算，奇数可先翻倍成偶数，之后只能不断把当前最大偶数减半；每次记录最大值和当前最小值差。

\textbf{重复工作。} 题面模型：奇数可以先乘二，偶数可以不断除二，目标是缩小当前最大最小差。暴力版的慢点要落到本题：慢版本会围绕“奇数可以先乘二，偶数可以不断除二，目标是缩小当前最大最小差”反复扫描动态候选集找极值。样例规律是“规律是把所有数变到可下降的最高状态，再用大顶堆逐步降低最大值”，说明每轮真正需要的是堆顶边界，而不是把候选全集重新排序。后面的优化只消掉这类重复工作，不能改变答案集合。

\textbf{优化条件。} 本题的关键观察是：把所有数规范到可降低的偶数形态，用大顶堆维护当前最大值，同时记录当前最小值。这句话必须能翻译成可维护状态；如果题目条件改变后观察不再成立，就不能继续套原来的写法。

\textbf{相邻状态。} 规律是把所有数变到可下降的最高状态，再用大顶堆逐步降低最大值。把这个特例继续拆开看：堆顶必须正好代表下一步最需要处理的候选；若堆里可能有过期对象，读取堆顶前要先清理。堆只解决动态极值，不自动证明被弹出或被丢弃的元素无用。

\textbf{状态复用。} 把每轮全量扫描改成堆顶查询；插入、弹出和过期清理共同维护动态候选集。本题落点是：规律是把所有数变到可下降的最高状态，再用大顶堆逐步降低最大值。手算时只改被新输入影响的那一小部分，而不是回头重算完整候选。

\textbf{指针和字段。} 状态字段要能说清：堆元素字段、堆顶比较规则、候选是否过期、答案读取时机；如果需要懒删除，还要保存删除计数。手算时按这个协议跟踪：拿 [1, 2, 3, 4] 手算，奇数可先翻倍成偶数，之后只能不断把当前最大偶数减半；每次记录最大值和当前最小值差。然后按协议复查：手算时记录堆顶、候选集合和是否有过期元素。每次使用堆顶前都要确认它仍然属于当前问题。

\textbf{代码落点。} 比较器方向先用小样例验证；每次使用堆顶前清理过期项；堆中保存下标时要能回查原数据。关键代码行必须能逐句翻译成“旧状态怎样变成新状态”，否则就只是模板。

\textbf{正确性。} 证明从候选覆盖开始：所有可能成为下一步答案的对象都在堆中，堆顶过期时不会被使用。

\textbf{边界实验。} 先用全奇数、全偶数、最大值变奇数停止、数值范围做反例检查。实验时覆盖全奇数、全偶数、最大值变奇数停止、数值范围，记录堆顶是否过期、比较器是否让正确候选在堆顶、K 或窗口大小为边界值时是否稳定。

\textbf{复杂度变化。} 暴力每轮扫描或排序动态候选，会把线性扫描或排序反复发生；堆把取极值、插入和删除边界控制在对数级。

\textbf{迁移追问。} 如果题目改成“这是堆维护动态极值和状态空间规范化的结合”，先判断原状态是否仍然覆盖新答案集合，再决定是微调边界、改 value 语义，还是换算法。

\begin{principlebox}{本节复盘}
阅读本节后，请回到相邻章节的 LeetCode 案例，例如 LeetCode 875 爱吃香蕉、LeetCode 72 编辑距离、LeetCode 743 网络延迟时间、LeetCode 215 数组第 K 大，把暴力瓶颈、优化依据、状态不变量、C++ 容器成本和边界测试重新写一遍。能从这些具体题迁移到变体题，才算真正掌握。
\end{principlebox}
